
\documentclass{article}
\usepackage{colortbl}
\usepackage{makecell}
\usepackage{multirow}
\usepackage{supertabular}

\begin{document}

\newcounter{utterance}

\centering \large Interaction Transcript for game `clean\_up', experiment `1\_medium\_7obj\_en', episode 1 with llama{-}3.3{-}70b{-}instruct{-}t1.0.
\vspace{24pt}

{ \footnotesize  \setcounter{utterance}{1}
\setlength{\tabcolsep}{0pt}
\begin{supertabular}{c@{$\;$}|p{.15\linewidth}@{}p{.15\linewidth}p{.15\linewidth}p{.15\linewidth}p{.15\linewidth}p{.15\linewidth}}
    \# & $\;$A & \multicolumn{4}{c}{Game Master} & $\;\:$B\\
    \hline

    \theutterance \stepcounter{utterance}  
    & & \multicolumn{4}{p{0.6\linewidth}}{
        \cellcolor[rgb]{0.9,0.9,0.9}{
            \makecell[{{p{\linewidth}}}]{
                \texttt{\tiny{[P1$\langle$GM]}}
                \texttt{Sono il tuo game master, e stai giocando un gioco collaborativo con la seguente griglia come tavolo da gioco:} \\
\\ 
\texttt{```} \\
\texttt{ 1234567} \\
\texttt{╔═══╤═╤═╗} \\
\texttt{║◌◌◌│◌│◌║ 1} \\
\texttt{║◌◌M│◌├─╢ 2} \\
\texttt{║AT◌├─┘◌║ 3} \\
\texttt{║◌◌P│◌◌◌║ 4} \\
\texttt{╟─┬─┼───╢ 5} \\
\texttt{║◌│◌│S◌O║ 6} \\
\texttt{║◌│◌│◌D◌║ 7} \\
\texttt{╚═╧═╧═══╝} \\
\\ 
\texttt{```} \\
\\ 
\texttt{* Il bordo superiore mostra le coordinate x che aumentano verso destra, e il bordo destro le coordinate y che aumentano verso il basso.} \\
\texttt{* The following objects are randomly placed on your grid: 'P', 'O', 'T', 'S', 'D', 'A', 'M'.} \\
\\ 
\\ 
\texttt{L'altro giocatore vede una variante del tavolo da gioco, dove gli oggetti sono posizionati in diverse posizioni casuali. Non puoi vedere il tavolo dell'altro giocatore, e loro non possono vedere il tuo.} \\
\\ 
\\ 
\texttt{**Obiettivo:**} \\
\\ 
\texttt{Entrambi i giocatori devono spostare gli oggetti sui rispettivi sfondi in modo che oggetti identici finiscano alle stesse coordinate. Devi comunicare con l'altro giocatore per concordare una configurazione obiettivo comune.} \\
\\ 
\\ 
\texttt{**Regole:**} \\
\\ 
\texttt{* In ogni turno, puoi inviare esattamente uno dei seguenti due comandi:} \\
\texttt{ 1. `COMUNICA: <MESSAGE>`: per inviare un messaggio (tutto fino alla prossima interruzione di riga) all'altro giocatore. Lo inoltrerò al tuo partner.} \\
\texttt{2. `SPOSTA: <OBJECT>, (<X>, <Y>)`: per spostare un oggetto in una nuova posizione, dove <X> è la colonna e <Y> è la riga. Ti farò sapere se la tua mossa è andata a buon fine o no.} \\
\texttt{* Se non rispetti il formato, o invii diversi comandi contemporaneamente, devo penalizzarti.} \\
\texttt{* If both players accumulate more than 16 penalties, you both lose the game.} \\
\texttt{* È vitale che comunichi con l'altro giocatore riguardo al tuo stato obiettivo! L'*unico* modo in cui puoi trasmettere la tua strategia all'altro giocatore è usando il comando `COMUNICA: <MESSAGE>`!} \\
\\ 
\\ 
\texttt{**Spostamento Oggetti**} \\
\\ 
\texttt{* Puoi spostare oggetti solo in celle entro i limiti della griglia. La cella di destinazione deve essere vuota, cioè deve contenere solo il simbolo '◌'.} \\
\texttt{* Se provi a spostare un oggetto in un posto che non è vuoto, o provi a spostarlo fuori dalla griglia, devo penalizzarti. Hai un altro tentativo.} \\
\texttt{* Prima di fare una mossa, controlla due volte che il posto di destinazione sia vuoto, e non contenga alcuna lettera, cornice, o linea!} \\
\\ 
\\ 
\texttt{**Fine del Gioco**} \\
\\ 
\texttt{Se pensi di aver raggiunto l'obiettivo di allineare tutti gli oggetti, puoi chiedere all'altro giocatore di finire il gioco inviando COMUNICA: finito?. Se l'altro giocatore ti chiede di finire il gioco, e rispondi con COMUNICA: finito!, the game will end.} \\
\\ 
\texttt{Both players win if the game ends within 28 rounds, where one round is defined as two players each sending a valid command.} \\
\\ 
\texttt{**Punteggio:**} \\
\\ 
\texttt{Più vicini sono gli oggetti identici in entrambi i tavoli da gioco, più punti ottieni. Le penalità riducono i tuoi punti. Riesci a battere il record?} \\
\\ 
\\ 
\texttt{Per favore invia un messaggio all'altro giocatore per iniziare il gioco!} \\
            }
        }
    }
    & \\ \\

    \theutterance \stepcounter{utterance}  
    & \multicolumn{4}{p{0.6\linewidth}}{
        \cellcolor[rgb]{0.9,0.9,0.9}{
            \makecell[{{p{\linewidth}}}]{
                \texttt{\tiny{[P1$\rangle$GM]}}
                \texttt{Sono emozionato di iniziare il gioco! } \\
\\ 
\texttt{COMUNICA: Ciao! Sono pronto a iniziare. La mia griglia attuale è:} \\
\\ 
\texttt{```} \\
\texttt{ 1234567} \\
\texttt{╔═══╤═╤═╗} \\
\texttt{║◌◌◌│◌│◌║ 1} \\
\texttt{║◌◌M│◌├─╢ 2} \\
\texttt{║AT◌├─┘◌║ 3} \\
\texttt{║◌◌P│◌◌◌║ 4} \\
\texttt{╟─┬─┼───╢ 5} \\
\texttt{║◌│◌│S◌O║ 6} \\
\texttt{║◌│◌│◌D◌║ 7} \\
\texttt{╚═╧═╧═══╝} \\
\texttt{```} \\
\\ 
\texttt{Propongo di iniziare a discutere sulla posizione degli oggetti. Qual è la tua strategia? Quali oggetti pensi di spostare per primo?} \\
            }
        }
    }
    & & \\ \\

    \theutterance \stepcounter{utterance}  
    & & & \multicolumn{2}{p{0.3\linewidth}}{
        \cellcolor[rgb]{0.9,0.9,0.9}{
            \makecell[{{p{\linewidth}}}]{
                \texttt{\tiny{[GM$|$GM]}}
                \texttt{Sono emozionato di iniziare il gioco! } \\
\\ 
\texttt{COMUNICA: Ciao! Sono pronto a iniziare. La mia griglia attuale è:} \\
\\ 
\texttt{```} \\
\texttt{ 1234567} \\
\texttt{╔═══╤═╤═╗} \\
\texttt{║◌◌◌│◌│◌║ 1} \\
\texttt{║◌◌M│◌├─╢ 2} \\
\texttt{║AT◌├─┘◌║ 3} \\
\texttt{║◌◌P│◌◌◌║ 4} \\
\texttt{╟─┬─┼───╢ 5} \\
\texttt{║◌│◌│S◌O║ 6} \\
\texttt{║◌│◌│◌D◌║ 7} \\
\texttt{╚═╧═╧═══╝} \\
\texttt{```} \\
\\ 
\texttt{Propongo di iniziare a discutere sulla posizione degli oggetti. Qual è la tua strategia? Quali oggetti pensi di spostare per primo?} \\
            }
        }
    }
    & & \\ \\

    \theutterance \stepcounter{utterance}  
    & & & \multicolumn{2}{p{0.3\linewidth}}{
        \cellcolor[rgb]{0.9,0.9,0.9}{
            \makecell[{{p{\linewidth}}}]{
                \texttt{\tiny{[GM$|$GM]}}
                \texttt{Invalid format: head and tail are not empty} \\
\texttt{head: 'Sono emozionato di iniziare il gioco! } \\
\\ 
\texttt{'} \\
\texttt{tail: '} \\
\\ 
\\ 
\texttt{ 1234567} \\
\texttt{╔═══╤═╤═╗} \\
\texttt{║◌◌◌│◌│◌║ 1} \\
\texttt{║◌◌M│◌├─╢ 2} \\
\texttt{║AT◌├─┘◌║ 3} \\
\texttt{║◌◌P│◌◌◌║ 4} \\
\texttt{╟─┬─┼───╢ 5} \\
\texttt{║◌│◌│S◌O║ 6} \\
\texttt{║◌│◌│◌D◌║ 7} \\
\texttt{╚═╧═╧═══╝} \\
\\ 
\\ 
\texttt{Propongo di iniziare a discutere sulla posizione degli oggetti. Qual è la tua strategia? Quali oggetti pensi di spostare per primo?'} \\
            }
        }
    }
    & & \\ \\

    \theutterance \stepcounter{utterance}  
    & & \multicolumn{4}{p{0.6\linewidth}}{
        \cellcolor[rgb]{0.9,0.9,0.9}{
            \makecell[{{p{\linewidth}}}]{
                \texttt{\tiny{[P1$\langle$GM]}}
                \texttt{Penalità: Il tuo messaggio non deve contenere nulla prima o dopo il comando!} \\
\\ 
\texttt{Assicurati che la tua risposta contenga solo `COMUNICA: <MESSAGE>` o `SPOSTA: <OBJECT>, (<X>, <Y>)`, e nient'altro!} \\
\\ 
\texttt{You have collectively accumulated 1 of 16 penalties.} \\
\texttt{Riprova!} \\
            }
        }
    }
    & \\ \\

    \theutterance \stepcounter{utterance}  
    & \multicolumn{4}{p{0.6\linewidth}}{
        \cellcolor[rgb]{0.9,0.9,0.9}{
            \makecell[{{p{\linewidth}}}]{
                \texttt{\tiny{[P1$\rangle$GM]}}
                \texttt{COMUNICA: Ciao, sono pronto a iniziare, la mia griglia attuale ha gli oggetti 'P', 'O', 'T', 'S', 'D', 'A', 'M' posizionati in varie celle, propongo di iniziare a discutere sulla posizione degli oggetti per trovare una configurazione comune.} \\
            }
        }
    }
    & & \\ \\

    \theutterance \stepcounter{utterance}  
    & & & \multicolumn{2}{p{0.3\linewidth}}{
        \cellcolor[rgb]{0.9,0.9,0.9}{
            \makecell[{{p{\linewidth}}}]{
                \texttt{\tiny{[GM$|$GM]}}
                \texttt{COMUNICA: Ciao, sono pronto a iniziare, la mia griglia attuale ha gli oggetti 'P', 'O', 'T', 'S', 'D', 'A', 'M' posizionati in varie celle, propongo di iniziare a discutere sulla posizione degli oggetti per trovare una configurazione comune.} \\
            }
        }
    }
    & & \\ \\

    \theutterance \stepcounter{utterance}  
    & & \multicolumn{4}{p{0.6\linewidth}}{
        \cellcolor[rgb]{0.9,0.9,0.9}{
            \makecell[{{p{\linewidth}}}]{
                \texttt{\tiny{[GM$\rangle$P2]}}
                \texttt{Sono il tuo game master, e stai giocando un gioco collaborativo con la seguente griglia come tavolo da gioco:} \\
\\ 
\texttt{```} \\
\texttt{ 1234567} \\
\texttt{╔═══╤═╤═╗} \\
\texttt{║◌MS│◌│◌║ 1} \\
\texttt{║◌O◌│◌├─╢ 2} \\
\texttt{║◌◌◌├─┘◌║ 3} \\
\texttt{║◌◌◌│◌◌◌║ 4} \\
\texttt{╟─┬─┼───╢ 5} \\
\texttt{║◌│◌│◌◌D║ 6} \\
\texttt{║T│◌│◌PA║ 7} \\
\texttt{╚═╧═╧═══╝} \\
\\ 
\texttt{```} \\
\\ 
\texttt{* Il bordo superiore mostra le coordinate x che aumentano verso destra, e il bordo destro le coordinate y che aumentano verso il basso.} \\
\texttt{* The following objects are randomly placed on your grid: 'P', 'O', 'T', 'S', 'D', 'A', 'M'.} \\
\\ 
\\ 
\texttt{L'altro giocatore vede una variante del tavolo da gioco, dove gli oggetti sono posizionati in diverse posizioni casuali. Non puoi vedere il tavolo dell'altro giocatore, e loro non possono vedere il tuo.} \\
\\ 
\\ 
\texttt{**Obiettivo:**} \\
\\ 
\texttt{Entrambi i giocatori devono spostare gli oggetti sui rispettivi sfondi in modo che oggetti identici finiscano alle stesse coordinate. Devi comunicare con l'altro giocatore per concordare una configurazione obiettivo comune.} \\
\\ 
\\ 
\texttt{**Regole:**} \\
\\ 
\texttt{* In ogni turno, puoi inviare esattamente uno dei seguenti due comandi:} \\
\texttt{ 1. `COMUNICA: <MESSAGE>`: per inviare un messaggio (tutto fino alla prossima interruzione di riga) all'altro giocatore. Lo inoltrerò al tuo partner.} \\
\texttt{2. `SPOSTA: <OBJECT>, (<X>, <Y>)`: per spostare un oggetto in una nuova posizione, dove <X> è la colonna e <Y> è la riga. Ti farò sapere se la tua mossa è andata a buon fine o no.} \\
\texttt{* Se non rispetti il formato, o invii diversi comandi contemporaneamente, devo penalizzarti.} \\
\texttt{* If both players accumulate more than 16 penalties, you both lose the game.} \\
\texttt{* È vitale che comunichi con l'altro giocatore riguardo al tuo stato obiettivo! L'*unico* modo in cui puoi trasmettere la tua strategia all'altro giocatore è usando il comando `COMUNICA: <MESSAGE>`!} \\
\\ 
\\ 
\texttt{**Spostamento Oggetti**} \\
\\ 
\texttt{* Puoi spostare oggetti solo in celle entro i limiti della griglia. La cella di destinazione deve essere vuota, cioè deve contenere solo il simbolo '◌'.} \\
\texttt{* Se provi a spostare un oggetto in un posto che non è vuoto, o provi a spostarlo fuori dalla griglia, devo penalizzarti. Hai un altro tentativo.} \\
\texttt{* Prima di fare una mossa, controlla due volte che il posto di destinazione sia vuoto, e non contenga alcuna lettera, cornice, o linea!} \\
\\ 
\\ 
\texttt{**Fine del Gioco**} \\
\\ 
\texttt{Se pensi di aver raggiunto l'obiettivo di allineare tutti gli oggetti, puoi chiedere all'altro giocatore di finire il gioco inviando COMUNICA: finito?. Se l'altro giocatore ti chiede di finire il gioco, e rispondi con COMUNICA: finito!, the game will end.} \\
\\ 
\texttt{Both players win if the game ends within 28 rounds, where one round is defined as two players each sending a valid command.} \\
\\ 
\texttt{**Punteggio:**} \\
\\ 
\texttt{Più vicini sono gli oggetti identici in entrambi i tavoli da gioco, più punti ottieni. Le penalità riducono i tuoi punti. Riesci a battere il record?} \\
\\ 
\texttt{L'altro giocatore ha iniziato il gioco inviando questo messaggio:} \\
\\ 
\texttt{"Ciao, sono pronto a iniziare, la mia griglia attuale ha gli oggetti 'P', 'O', 'T', 'S', 'D', 'A', 'M' posizionati in varie celle, propongo di iniziare a discutere sulla posizione degli oggetti per trovare una configurazione comune."} \\
\\ 
\texttt{Qual è il tuo primo comando?} \\
            }
        }
    }
    & \\ \\

    \theutterance \stepcounter{utterance}  
    & & & \multicolumn{4}{p{0.6\linewidth}}{
        \cellcolor[rgb]{0.9,0.9,0.9}{
            \makecell[{{p{\linewidth}}}]{
                \texttt{\tiny{[GM$\langle$P2]}}
                \texttt{COMUNICA: Grazie per il messaggio! Sono anche io pronto a iniziare. La mia griglia ha gli oggetti 'P', 'O', 'T', 'S', 'D', 'A', 'M' posizionati in varie celle, ma non so ancora quali siano le posizioni esatte. Propongo di iniziare a discutere sulle posizioni degli oggetti uno per uno, partendo da 'P'. Qual è la posizione del tuo 'P'?} \\
            }
        }
    }
     \\ \\

    \theutterance \stepcounter{utterance}  
    & & & \multicolumn{2}{p{0.3\linewidth}}{
        \cellcolor[rgb]{0.9,0.9,0.9}{
            \makecell[{{p{\linewidth}}}]{
                \texttt{\tiny{[GM$|$GM]}}
                \texttt{COMUNICA: Grazie per il messaggio! Sono anche io pronto a iniziare. La mia griglia ha gli oggetti 'P', 'O', 'T', 'S', 'D', 'A', 'M' posizionati in varie celle, ma non so ancora quali siano le posizioni esatte. Propongo di iniziare a discutere sulle posizioni degli oggetti uno per uno, partendo da 'P'. Qual è la posizione del tuo 'P'?} \\
            }
        }
    }
    & & \\ \\

    \theutterance \stepcounter{utterance}  
    & & \multicolumn{4}{p{0.6\linewidth}}{
        \cellcolor[rgb]{0.9,0.9,0.9}{
            \makecell[{{p{\linewidth}}}]{
                \texttt{\tiny{[P1$\langle$GM]}}
                \texttt{Il tuo messaggio è stato inoltrato all'altro giocatore.} \\
\\ 
\\ 
\texttt{You are currently playing round 1 of maximum 28.} \\
\texttt{You have collectively accumulated 1 of 16 penalties.} \\
\texttt{L'altro giocatore ha inviato questo messaggio:} \\
\\ 
\\ 
\texttt{"Grazie per il messaggio! Sono anche io pronto a iniziare. La mia griglia ha gli oggetti 'P', 'O', 'T', 'S', 'D', 'A', 'M' posizionati in varie celle, ma non so ancora quali siano le posizioni esatte. Propongo di iniziare a discutere sulle posizioni degli oggetti uno per uno, partendo da 'P'. Qual è la posizione del tuo 'P'?"} \\
\\ 
\\ 
\texttt{Qual è il tuo prossimo comando?} \\
            }
        }
    }
    & \\ \\

    \theutterance \stepcounter{utterance}  
    & \multicolumn{4}{p{0.6\linewidth}}{
        \cellcolor[rgb]{0.9,0.9,0.9}{
            \makecell[{{p{\linewidth}}}]{
                \texttt{\tiny{[P1$\rangle$GM]}}
                \texttt{COMUNICA: Il mio 'P' è posizionato alla riga 4, penso di spostarlo alla riga 1, colonna 1, ma prima vorrei sapere quale posizione del 'P' hai tu nella tua griglia.} \\
            }
        }
    }
    & & \\ \\

    \theutterance \stepcounter{utterance}  
    & & & \multicolumn{2}{p{0.3\linewidth}}{
        \cellcolor[rgb]{0.9,0.9,0.9}{
            \makecell[{{p{\linewidth}}}]{
                \texttt{\tiny{[GM$|$GM]}}
                \texttt{COMUNICA: Il mio 'P' è posizionato alla riga 4, penso di spostarlo alla riga 1, colonna 1, ma prima vorrei sapere quale posizione del 'P' hai tu nella tua griglia.} \\
            }
        }
    }
    & & \\ \\

    \theutterance \stepcounter{utterance}  
    & & \multicolumn{4}{p{0.6\linewidth}}{
        \cellcolor[rgb]{0.9,0.9,0.9}{
            \makecell[{{p{\linewidth}}}]{
                \texttt{\tiny{[GM$\rangle$P2]}}
                \texttt{Il tuo messaggio è stato inoltrato all'altro giocatore.} \\
\\ 
\\ 
\texttt{You are currently playing round 2 of maximum 28.} \\
\texttt{You have collectively accumulated 1 of 16 penalties.} \\
\texttt{L'altro giocatore ha inviato questo messaggio:} \\
\\ 
\\ 
\texttt{"Il mio 'P' è posizionato alla riga 4, penso di spostarlo alla riga 1, colonna 1, ma prima vorrei sapere quale posizione del 'P' hai tu nella tua griglia."} \\
\\ 
\\ 
\texttt{Qual è il tuo prossimo comando?} \\
            }
        }
    }
    & \\ \\

    \theutterance \stepcounter{utterance}  
    & & & \multicolumn{4}{p{0.6\linewidth}}{
        \cellcolor[rgb]{0.9,0.9,0.9}{
            \makecell[{{p{\linewidth}}}]{
                \texttt{\tiny{[GM$\langle$P2]}}
                \texttt{COMUNICA: Il mio 'P' è attualmente alla riga 7, colonna 3, nella mia griglia. Considerando che tu vuoi spostare il tuo 'P' alla riga 1, colonna 1, penso che potremmo raggiungere una configurazione comune spostando entrambi i nostri 'P' nella stessa posizione. Tuttavia, prima di procedere, vorrei sapere qualcosa sugli altri oggetti. Qual è la posizione del tuo 'T'?} \\
            }
        }
    }
     \\ \\

    \theutterance \stepcounter{utterance}  
    & & & \multicolumn{2}{p{0.3\linewidth}}{
        \cellcolor[rgb]{0.9,0.9,0.9}{
            \makecell[{{p{\linewidth}}}]{
                \texttt{\tiny{[GM$|$GM]}}
                \texttt{COMUNICA: Il mio 'P' è attualmente alla riga 7, colonna 3, nella mia griglia. Considerando che tu vuoi spostare il tuo 'P' alla riga 1, colonna 1, penso che potremmo raggiungere una configurazione comune spostando entrambi i nostri 'P' nella stessa posizione. Tuttavia, prima di procedere, vorrei sapere qualcosa sugli altri oggetti. Qual è la posizione del tuo 'T'?} \\
            }
        }
    }
    & & \\ \\

    \theutterance \stepcounter{utterance}  
    & & \multicolumn{4}{p{0.6\linewidth}}{
        \cellcolor[rgb]{0.9,0.9,0.9}{
            \makecell[{{p{\linewidth}}}]{
                \texttt{\tiny{[P1$\langle$GM]}}
                \texttt{Il tuo messaggio è stato inoltrato all'altro giocatore.} \\
\\ 
\\ 
\texttt{You are currently playing round 2 of maximum 28.} \\
\texttt{You have collectively accumulated 1 of 16 penalties.} \\
\texttt{L'altro giocatore ha inviato questo messaggio:} \\
\\ 
\\ 
\texttt{"Il mio 'P' è attualmente alla riga 7, colonna 3, nella mia griglia. Considerando che tu vuoi spostare il tuo 'P' alla riga 1, colonna 1, penso che potremmo raggiungere una configurazione comune spostando entrambi i nostri 'P' nella stessa posizione. Tuttavia, prima di procedere, vorrei sapere qualcosa sugli altri oggetti. Qual è la posizione del tuo 'T'?"} \\
\\ 
\\ 
\texttt{Qual è il tuo prossimo comando?} \\
            }
        }
    }
    & \\ \\

    \theutterance \stepcounter{utterance}  
    & \multicolumn{4}{p{0.6\linewidth}}{
        \cellcolor[rgb]{0.9,0.9,0.9}{
            \makecell[{{p{\linewidth}}}]{
                \texttt{\tiny{[P1$\rangle$GM]}}
                \texttt{COMUNICA: Il mio 'T' è attualmente alla riga 3, ed è adiacente ad 'A', mentre il tuo 'P' è alla riga 7, colonna 3, nella tua griglia, quindi potremmo avere una diversa disposizione, forse potremmo concentrarci su 'S' e 'D' per trovare una posizione comune per loro.} \\
            }
        }
    }
    & & \\ \\

    \theutterance \stepcounter{utterance}  
    & & & \multicolumn{2}{p{0.3\linewidth}}{
        \cellcolor[rgb]{0.9,0.9,0.9}{
            \makecell[{{p{\linewidth}}}]{
                \texttt{\tiny{[GM$|$GM]}}
                \texttt{COMUNICA: Il mio 'T' è attualmente alla riga 3, ed è adiacente ad 'A', mentre il tuo 'P' è alla riga 7, colonna 3, nella tua griglia, quindi potremmo avere una diversa disposizione, forse potremmo concentrarci su 'S' e 'D' per trovare una posizione comune per loro.} \\
            }
        }
    }
    & & \\ \\

    \theutterance \stepcounter{utterance}  
    & & \multicolumn{4}{p{0.6\linewidth}}{
        \cellcolor[rgb]{0.9,0.9,0.9}{
            \makecell[{{p{\linewidth}}}]{
                \texttt{\tiny{[GM$\rangle$P2]}}
                \texttt{Il tuo messaggio è stato inoltrato all'altro giocatore.} \\
\\ 
\\ 
\texttt{You are currently playing round 3 of maximum 28.} \\
\texttt{You have collectively accumulated 1 of 16 penalties.} \\
\texttt{L'altro giocatore ha inviato questo messaggio:} \\
\\ 
\\ 
\texttt{"Il mio 'T' è attualmente alla riga 3, ed è adiacente ad 'A', mentre il tuo 'P' è alla riga 7, colonna 3, nella tua griglia, quindi potremmo avere una diversa disposizione, forse potremmo concentrarci su 'S' e 'D' per trovare una posizione comune per loro."} \\
\\ 
\\ 
\texttt{Qual è il tuo prossimo comando?} \\
            }
        }
    }
    & \\ \\

    \theutterance \stepcounter{utterance}  
    & & & \multicolumn{4}{p{0.6\linewidth}}{
        \cellcolor[rgb]{0.9,0.9,0.9}{
            \makecell[{{p{\linewidth}}}]{
                \texttt{\tiny{[GM$\langle$P2]}}
                \texttt{COMUNICA: Sì, concentrarsi su 'S' e 'D' potrebbe essere una buona idea. Nella mia griglia, 'S' è alla riga 1, colonna 2, e 'D' è alla riga 6, colonna 4. Come sono disposti 'S' e 'D' nella tua griglia? Inoltre, pensi che potremmo trovare una posizione centrale per 'S' e 'D' in modo che sia facile allineare gli altri oggetti in seguito?} \\
            }
        }
    }
     \\ \\

    \theutterance \stepcounter{utterance}  
    & & & \multicolumn{2}{p{0.3\linewidth}}{
        \cellcolor[rgb]{0.9,0.9,0.9}{
            \makecell[{{p{\linewidth}}}]{
                \texttt{\tiny{[GM$|$GM]}}
                \texttt{COMUNICA: Sì, concentrarsi su 'S' e 'D' potrebbe essere una buona idea. Nella mia griglia, 'S' è alla riga 1, colonna 2, e 'D' è alla riga 6, colonna 4. Come sono disposti 'S' e 'D' nella tua griglia? Inoltre, pensi che potremmo trovare una posizione centrale per 'S' e 'D' in modo che sia facile allineare gli altri oggetti in seguito?} \\
            }
        }
    }
    & & \\ \\

    \theutterance \stepcounter{utterance}  
    & & \multicolumn{4}{p{0.6\linewidth}}{
        \cellcolor[rgb]{0.9,0.9,0.9}{
            \makecell[{{p{\linewidth}}}]{
                \texttt{\tiny{[P1$\langle$GM]}}
                \texttt{Il tuo messaggio è stato inoltrato all'altro giocatore.} \\
\\ 
\\ 
\texttt{You are currently playing round 3 of maximum 28.} \\
\texttt{You have collectively accumulated 1 of 16 penalties.} \\
\texttt{L'altro giocatore ha inviato questo messaggio:} \\
\\ 
\\ 
\texttt{"Sì, concentrarsi su 'S' e 'D' potrebbe essere una buona idea. Nella mia griglia, 'S' è alla riga 1, colonna 2, e 'D' è alla riga 6, colonna 4. Come sono disposti 'S' e 'D' nella tua griglia? Inoltre, pensi che potremmo trovare una posizione centrale per 'S' e 'D' in modo che sia facile allineare gli altri oggetti in seguito?"} \\
\\ 
\\ 
\texttt{Qual è il tuo prossimo comando?} \\
            }
        }
    }
    & \\ \\

    \theutterance \stepcounter{utterance}  
    & \multicolumn{4}{p{0.6\linewidth}}{
        \cellcolor[rgb]{0.9,0.9,0.9}{
            \makecell[{{p{\linewidth}}}]{
                \texttt{\tiny{[P1$\rangle$GM]}}
                \texttt{COMUNICA: Nella mia griglia, 'S' è alla riga 6, colonna 3, e 'D' è alla riga 7, colonna 4, quindi abbiamo una disposizione diversa, ma possiamo lavorare per trovare una posizione centrale per 'S' e 'D', forse la riga 4 o 5 potrebbe essere un buon punto di incontro, ma dobbiamo anche considerare la posizione di 'M' e 'A' per non crearci problemi futuri.} \\
            }
        }
    }
    & & \\ \\

    \theutterance \stepcounter{utterance}  
    & & & \multicolumn{2}{p{0.3\linewidth}}{
        \cellcolor[rgb]{0.9,0.9,0.9}{
            \makecell[{{p{\linewidth}}}]{
                \texttt{\tiny{[GM$|$GM]}}
                \texttt{COMUNICA: Nella mia griglia, 'S' è alla riga 6, colonna 3, e 'D' è alla riga 7, colonna 4, quindi abbiamo una disposizione diversa, ma possiamo lavorare per trovare una posizione centrale per 'S' e 'D', forse la riga 4 o 5 potrebbe essere un buon punto di incontro, ma dobbiamo anche considerare la posizione di 'M' e 'A' per non crearci problemi futuri.} \\
            }
        }
    }
    & & \\ \\

    \theutterance \stepcounter{utterance}  
    & & \multicolumn{4}{p{0.6\linewidth}}{
        \cellcolor[rgb]{0.9,0.9,0.9}{
            \makecell[{{p{\linewidth}}}]{
                \texttt{\tiny{[GM$\rangle$P2]}}
                \texttt{Il tuo messaggio è stato inoltrato all'altro giocatore.} \\
\\ 
\\ 
\texttt{You are currently playing round 4 of maximum 28.} \\
\texttt{You have collectively accumulated 1 of 16 penalties.} \\
\texttt{L'altro giocatore ha inviato questo messaggio:} \\
\\ 
\\ 
\texttt{"Nella mia griglia, 'S' è alla riga 6, colonna 3, e 'D' è alla riga 7, colonna 4, quindi abbiamo una disposizione diversa, ma possiamo lavorare per trovare una posizione centrale per 'S' e 'D', forse la riga 4 o 5 potrebbe essere un buon punto di incontro, ma dobbiamo anche considerare la posizione di 'M' e 'A' per non crearci problemi futuri."} \\
\\ 
\\ 
\texttt{Qual è il tuo prossimo comando?} \\
            }
        }
    }
    & \\ \\

    \theutterance \stepcounter{utterance}  
    & & & \multicolumn{4}{p{0.6\linewidth}}{
        \cellcolor[rgb]{0.9,0.9,0.9}{
            \makecell[{{p{\linewidth}}}]{
                \texttt{\tiny{[GM$\langle$P2]}}
                \texttt{COMUNICA: Sì, la riga 4 o 5 potrebbe essere un buon punto di incontro per 'S' e 'D'. Nella mia griglia, 'M' è alla riga 1, colonna 1, e 'A' è alla riga 7, colonna 3, vicino a 'P'. Considerando le posizioni che hai descritto, potremmo cercare di spostare 'S' e 'D' verso la riga 5, in modo da lasciare spazio per 'M' e 'A' nelle rispettive griglie. Qual è la tua idea su come potremmo spostare 'M' e 'A' per non creare problemi futuri? Inoltre, vorrei sapere, hai già una strategia per spostare 'O' nella tua griglia?} \\
            }
        }
    }
     \\ \\

    \theutterance \stepcounter{utterance}  
    & & & \multicolumn{2}{p{0.3\linewidth}}{
        \cellcolor[rgb]{0.9,0.9,0.9}{
            \makecell[{{p{\linewidth}}}]{
                \texttt{\tiny{[GM$|$GM]}}
                \texttt{COMUNICA: Sì, la riga 4 o 5 potrebbe essere un buon punto di incontro per 'S' e 'D'. Nella mia griglia, 'M' è alla riga 1, colonna 1, e 'A' è alla riga 7, colonna 3, vicino a 'P'. Considerando le posizioni che hai descritto, potremmo cercare di spostare 'S' e 'D' verso la riga 5, in modo da lasciare spazio per 'M' e 'A' nelle rispettive griglie. Qual è la tua idea su come potremmo spostare 'M' e 'A' per non creare problemi futuri? Inoltre, vorrei sapere, hai già una strategia per spostare 'O' nella tua griglia?} \\
            }
        }
    }
    & & \\ \\

    \theutterance \stepcounter{utterance}  
    & & \multicolumn{4}{p{0.6\linewidth}}{
        \cellcolor[rgb]{0.9,0.9,0.9}{
            \makecell[{{p{\linewidth}}}]{
                \texttt{\tiny{[P1$\langle$GM]}}
                \texttt{Il tuo messaggio è stato inoltrato all'altro giocatore.} \\
\\ 
\\ 
\texttt{You are currently playing round 4 of maximum 28.} \\
\texttt{You have collectively accumulated 1 of 16 penalties.} \\
\texttt{L'altro giocatore ha inviato questo messaggio:} \\
\\ 
\\ 
\texttt{"Sì, la riga 4 o 5 potrebbe essere un buon punto di incontro per 'S' e 'D'. Nella mia griglia, 'M' è alla riga 1, colonna 1, e 'A' è alla riga 7, colonna 3, vicino a 'P'. Considerando le posizioni che hai descritto, potremmo cercare di spostare 'S' e 'D' verso la riga 5, in modo da lasciare spazio per 'M' e 'A' nelle rispettive griglie. Qual è la tua idea su come potremmo spostare 'M' e 'A' per non creare problemi futuri? Inoltre, vorrei sapere, hai già una strategia per spostare 'O' nella tua griglia?"} \\
\\ 
\\ 
\texttt{Qual è il tuo prossimo comando?} \\
            }
        }
    }
    & \\ \\

    \theutterance \stepcounter{utterance}  
    & \multicolumn{4}{p{0.6\linewidth}}{
        \cellcolor[rgb]{0.9,0.9,0.9}{
            \makecell[{{p{\linewidth}}}]{
                \texttt{\tiny{[P1$\rangle$GM]}}
                \texttt{COMUNICA: Nella mia griglia, 'M' è alla riga 2, colonna 3, e 'A' è alla riga 3, colonna 1, vicino a 'T', quindi abbiamo una disposizione un po' diversa, per 'O' penso di spostarlo vicino a 'S' e 'D' una volta che troveremo una posizione centrale per loro, magari alla riga 5, in questo modo potremmo creare un cluster di oggetti e facilitare il processo di allineamento, tu cosa ne pensi?} \\
            }
        }
    }
    & & \\ \\

    \theutterance \stepcounter{utterance}  
    & & & \multicolumn{2}{p{0.3\linewidth}}{
        \cellcolor[rgb]{0.9,0.9,0.9}{
            \makecell[{{p{\linewidth}}}]{
                \texttt{\tiny{[GM$|$GM]}}
                \texttt{COMUNICA: Nella mia griglia, 'M' è alla riga 2, colonna 3, e 'A' è alla riga 3, colonna 1, vicino a 'T', quindi abbiamo una disposizione un po' diversa, per 'O' penso di spostarlo vicino a 'S' e 'D' una volta che troveremo una posizione centrale per loro, magari alla riga 5, in questo modo potremmo creare un cluster di oggetti e facilitare il processo di allineamento, tu cosa ne pensi?} \\
            }
        }
    }
    & & \\ \\

    \theutterance \stepcounter{utterance}  
    & & \multicolumn{4}{p{0.6\linewidth}}{
        \cellcolor[rgb]{0.9,0.9,0.9}{
            \makecell[{{p{\linewidth}}}]{
                \texttt{\tiny{[GM$\rangle$P2]}}
                \texttt{Il tuo messaggio è stato inoltrato all'altro giocatore.} \\
\\ 
\\ 
\texttt{You are currently playing round 5 of maximum 28.} \\
\texttt{You have collectively accumulated 1 of 16 penalties.} \\
\texttt{L'altro giocatore ha inviato questo messaggio:} \\
\\ 
\\ 
\texttt{"Nella mia griglia, 'M' è alla riga 2, colonna 3, e 'A' è alla riga 3, colonna 1, vicino a 'T', quindi abbiamo una disposizione un po' diversa, per 'O' penso di spostarlo vicino a 'S' e 'D' una volta che troveremo una posizione centrale per loro, magari alla riga 5, in questo modo potremmo creare un cluster di oggetti e facilitare il processo di allineamento, tu cosa ne pensi?"} \\
\\ 
\\ 
\texttt{Qual è il tuo prossimo comando?} \\
            }
        }
    }
    & \\ \\

    \theutterance \stepcounter{utterance}  
    & & & \multicolumn{4}{p{0.6\linewidth}}{
        \cellcolor[rgb]{0.9,0.9,0.9}{
            \makecell[{{p{\linewidth}}}]{
                \texttt{\tiny{[GM$\langle$P2]}}
                \texttt{COMUNICA: Sì, creare un cluster di oggetti vicino alla riga 5 potrebbe essere una buona strategia per facilitare l'allineamento. Nella mia griglia, 'O' è attualmente alla riga 2, colonna 2, quindi potrei spostarlo più in basso una volta che avremo trovato una posizione centrale per 'S' e 'D'. Mi sembra che stiamo arrivando a una visione condivisa, quindi propongo di iniziare a fare Some mosse concrete per allineare gli oggetti. Per prima cosa, vorrei spostare 'P' nella mia griglia alla riga 1, colonna 1, per allinearlo con la posizione che hai descritto in precedenza. Ti va di fare la stessa mossa con il tuo 'P'?} \\
            }
        }
    }
     \\ \\

    \theutterance \stepcounter{utterance}  
    & & & \multicolumn{2}{p{0.3\linewidth}}{
        \cellcolor[rgb]{0.9,0.9,0.9}{
            \makecell[{{p{\linewidth}}}]{
                \texttt{\tiny{[GM$|$GM]}}
                \texttt{COMUNICA: Sì, creare un cluster di oggetti vicino alla riga 5 potrebbe essere una buona strategia per facilitare l'allineamento. Nella mia griglia, 'O' è attualmente alla riga 2, colonna 2, quindi potrei spostarlo più in basso una volta che avremo trovato una posizione centrale per 'S' e 'D'. Mi sembra che stiamo arrivando a una visione condivisa, quindi propongo di iniziare a fare Some mosse concrete per allineare gli oggetti. Per prima cosa, vorrei spostare 'P' nella mia griglia alla riga 1, colonna 1, per allinearlo con la posizione che hai descritto in precedenza. Ti va di fare la stessa mossa con il tuo 'P'?} \\
            }
        }
    }
    & & \\ \\

    \theutterance \stepcounter{utterance}  
    & & \multicolumn{4}{p{0.6\linewidth}}{
        \cellcolor[rgb]{0.9,0.9,0.9}{
            \makecell[{{p{\linewidth}}}]{
                \texttt{\tiny{[P1$\langle$GM]}}
                \texttt{Il tuo messaggio è stato inoltrato all'altro giocatore.} \\
\\ 
\\ 
\texttt{You are currently playing round 5 of maximum 28.} \\
\texttt{You have collectively accumulated 1 of 16 penalties.} \\
\texttt{L'altro giocatore ha inviato questo messaggio:} \\
\\ 
\\ 
\texttt{"Sì, creare un cluster di oggetti vicino alla riga 5 potrebbe essere una buona strategia per facilitare l'allineamento. Nella mia griglia, 'O' è attualmente alla riga 2, colonna 2, quindi potrei spostarlo più in basso una volta che avremo trovato una posizione centrale per 'S' e 'D'. Mi sembra che stiamo arrivando a una visione condivisa, quindi propongo di iniziare a fare Some mosse concrete per allineare gli oggetti. Per prima cosa, vorrei spostare 'P' nella mia griglia alla riga 1, colonna 1, per allinearlo con la posizione che hai descritto in precedenza. Ti va di fare la stessa mossa con il tuo 'P'?"} \\
\\ 
\\ 
\texttt{Qual è il tuo prossimo comando?} \\
            }
        }
    }
    & \\ \\

    \theutterance \stepcounter{utterance}  
    & \multicolumn{4}{p{0.6\linewidth}}{
        \cellcolor[rgb]{0.9,0.9,0.9}{
            \makecell[{{p{\linewidth}}}]{
                \texttt{\tiny{[P1$\rangle$GM]}}
                \texttt{SPOSTA: P, (1, 1)} \\
            }
        }
    }
    & & \\ \\

    \theutterance \stepcounter{utterance}  
    & & & \multicolumn{2}{p{0.3\linewidth}}{
        \cellcolor[rgb]{0.9,0.9,0.9}{
            \makecell[{{p{\linewidth}}}]{
                \texttt{\tiny{[GM$|$GM]}}
                \texttt{SPOSTA: P, (1, 1)} \\
            }
        }
    }
    & & \\ \\

    \theutterance \stepcounter{utterance}  
    & & & \multicolumn{2}{p{0.3\linewidth}}{
        \cellcolor[rgb]{0.9,0.9,0.9}{
            \makecell[{{p{\linewidth}}}]{
                \texttt{\tiny{[GM$|$GM]}}
                \texttt{Spostato 'P' in (1,1) con successo. La tua griglia aggiornata appare così:} \\
\texttt{```} \\
\texttt{ 1234567} \\
\texttt{╔═══╤═╤═╗} \\
\texttt{║P◌◌│◌│◌║ 1} \\
\texttt{║◌◌M│◌├─╢ 2} \\
\texttt{║AT◌├─┘◌║ 3} \\
\texttt{║◌◌◌│◌◌◌║ 4} \\
\texttt{╟─┬─┼───╢ 5} \\
\texttt{║◌│◌│S◌O║ 6} \\
\texttt{║◌│◌│◌D◌║ 7} \\
\texttt{╚═╧═╧═══╝} \\
\texttt{```} \\
            }
        }
    }
    & & \\ \\

    \theutterance \stepcounter{utterance}  
    & & \multicolumn{4}{p{0.6\linewidth}}{
        \cellcolor[rgb]{0.9,0.9,0.9}{
            \makecell[{{p{\linewidth}}}]{
                \texttt{\tiny{[GM$\rangle$P2]}}
                \texttt{Il tuo messaggio è stato inoltrato all'altro giocatore.} \\
\\ 
\\ 
\texttt{You are currently playing round 6 of maximum 28.} \\
\texttt{You have collectively accumulated 1 of 16 penalties.} \\
\texttt{L'altro giocatore ha spostato un oggetto sulla sua griglia.} \\
\\ 
\\ 
\texttt{Qual è il tuo prossimo comando?} \\
            }
        }
    }
    & \\ \\

    \theutterance \stepcounter{utterance}  
    & & & \multicolumn{4}{p{0.6\linewidth}}{
        \cellcolor[rgb]{0.9,0.9,0.9}{
            \makecell[{{p{\linewidth}}}]{
                \texttt{\tiny{[GM$\langle$P2]}}
                \texttt{SPOSTA: P, (1, 1)} \\
            }
        }
    }
     \\ \\

    \theutterance \stepcounter{utterance}  
    & & & \multicolumn{2}{p{0.3\linewidth}}{
        \cellcolor[rgb]{0.9,0.9,0.9}{
            \makecell[{{p{\linewidth}}}]{
                \texttt{\tiny{[GM$|$GM]}}
                \texttt{SPOSTA: P, (1, 1)} \\
            }
        }
    }
    & & \\ \\

    \theutterance \stepcounter{utterance}  
    & & & \multicolumn{2}{p{0.3\linewidth}}{
        \cellcolor[rgb]{0.9,0.9,0.9}{
            \makecell[{{p{\linewidth}}}]{
                \texttt{\tiny{[GM$|$GM]}}
                \texttt{Spostato 'P' in (1,1) con successo. La tua griglia aggiornata appare così:} \\
\texttt{```} \\
\texttt{ 1234567} \\
\texttt{╔═══╤═╤═╗} \\
\texttt{║PMS│◌│◌║ 1} \\
\texttt{║◌O◌│◌├─╢ 2} \\
\texttt{║◌◌◌├─┘◌║ 3} \\
\texttt{║◌◌◌│◌◌◌║ 4} \\
\texttt{╟─┬─┼───╢ 5} \\
\texttt{║◌│◌│◌◌D║ 6} \\
\texttt{║T│◌│◌◌A║ 7} \\
\texttt{╚═╧═╧═══╝} \\
\texttt{```} \\
            }
        }
    }
    & & \\ \\

    \theutterance \stepcounter{utterance}  
    & & \multicolumn{4}{p{0.6\linewidth}}{
        \cellcolor[rgb]{0.9,0.9,0.9}{
            \makecell[{{p{\linewidth}}}]{
                \texttt{\tiny{[P1$\langle$GM]}}
                \texttt{Spostato 'P' in (1,1) con successo. La tua griglia aggiornata appare così:} \\
\texttt{```} \\
\texttt{ 1234567} \\
\texttt{╔═══╤═╤═╗} \\
\texttt{║P◌◌│◌│◌║ 1} \\
\texttt{║◌◌M│◌├─╢ 2} \\
\texttt{║AT◌├─┘◌║ 3} \\
\texttt{║◌◌◌│◌◌◌║ 4} \\
\texttt{╟─┬─┼───╢ 5} \\
\texttt{║◌│◌│S◌O║ 6} \\
\texttt{║◌│◌│◌D◌║ 7} \\
\texttt{╚═╧═╧═══╝} \\
\texttt{```} \\
\texttt{You are currently playing round 6 of maximum 28.} \\
\texttt{You have collectively accumulated 1 of 16 penalties.} \\
\texttt{L'altro giocatore ha spostato un oggetto sulla sua griglia.} \\
\\ 
\\ 
\texttt{Qual è il tuo prossimo comando?} \\
            }
        }
    }
    & \\ \\

    \theutterance \stepcounter{utterance}  
    & \multicolumn{4}{p{0.6\linewidth}}{
        \cellcolor[rgb]{0.9,0.9,0.9}{
            \makecell[{{p{\linewidth}}}]{
                \texttt{\tiny{[P1$\rangle$GM]}}
                \texttt{COMUNICA: Ho spostato il mio 'P' in (1,1), adesso mi concentro su 'S' e 'D', forse potremmo spostare 'S' alla riga 5, come discusso in precedenza, per allinearli meglio, tu cosa hai fatto con il tuo 'P' e come stanno 'S' e 'D' nella tua griglia adesso?} \\
            }
        }
    }
    & & \\ \\

    \theutterance \stepcounter{utterance}  
    & & & \multicolumn{2}{p{0.3\linewidth}}{
        \cellcolor[rgb]{0.9,0.9,0.9}{
            \makecell[{{p{\linewidth}}}]{
                \texttt{\tiny{[GM$|$GM]}}
                \texttt{COMUNICA: Ho spostato il mio 'P' in (1,1), adesso mi concentro su 'S' e 'D', forse potremmo spostare 'S' alla riga 5, come discusso in precedenza, per allinearli meglio, tu cosa hai fatto con il tuo 'P' e come stanno 'S' e 'D' nella tua griglia adesso?} \\
            }
        }
    }
    & & \\ \\

    \theutterance \stepcounter{utterance}  
    & & \multicolumn{4}{p{0.6\linewidth}}{
        \cellcolor[rgb]{0.9,0.9,0.9}{
            \makecell[{{p{\linewidth}}}]{
                \texttt{\tiny{[GM$\rangle$P2]}}
                \texttt{Spostato 'P' in (1,1) con successo. La tua griglia aggiornata appare così:} \\
\texttt{```} \\
\texttt{ 1234567} \\
\texttt{╔═══╤═╤═╗} \\
\texttt{║PMS│◌│◌║ 1} \\
\texttt{║◌O◌│◌├─╢ 2} \\
\texttt{║◌◌◌├─┘◌║ 3} \\
\texttt{║◌◌◌│◌◌◌║ 4} \\
\texttt{╟─┬─┼───╢ 5} \\
\texttt{║◌│◌│◌◌D║ 6} \\
\texttt{║T│◌│◌◌A║ 7} \\
\texttt{╚═╧═╧═══╝} \\
\texttt{```} \\
\texttt{You are currently playing round 7 of maximum 28.} \\
\texttt{You have collectively accumulated 1 of 16 penalties.} \\
\texttt{L'altro giocatore ha inviato questo messaggio:} \\
\\ 
\\ 
\texttt{"Ho spostato il mio 'P' in (1,1), adesso mi concentro su 'S' e 'D', forse potremmo spostare 'S' alla riga 5, come discusso in precedenza, per allinearli meglio, tu cosa hai fatto con il tuo 'P' e come stanno 'S' e 'D' nella tua griglia adesso?"} \\
\\ 
\\ 
\texttt{Qual è il tuo prossimo comando?} \\
            }
        }
    }
    & \\ \\

    \theutterance \stepcounter{utterance}  
    & & & \multicolumn{4}{p{0.6\linewidth}}{
        \cellcolor[rgb]{0.9,0.9,0.9}{
            \makecell[{{p{\linewidth}}}]{
                \texttt{\tiny{[GM$\langle$P2]}}
                \texttt{COMUNICA: Ho spostato il mio 'P' in (1,1) come concordato. Per quanto riguarda 'S' e 'D', nella mia griglia, 'S' è ancora alla riga 1, colonna 2, e 'D' è alla riga 6, colonna 4. Spostare 'S' alla riga 5 potrebbe essere una buona idea, ma dobbiamo anche considerare la posizione di 'M' e altri oggetti per non creare conflitti. Prima di spostare 'S', vorrei sapere come stanno 'M' e 'A' nella tua griglia adesso, e se ci sono altre mosse che puoi fare per allineare gli oggetti. Inoltre, potresti dirmi se hai già pensato a come spostare 'O' vicino a 'S' e 'D' una volta che li avremo allineati?} \\
            }
        }
    }
     \\ \\

    \theutterance \stepcounter{utterance}  
    & & & \multicolumn{2}{p{0.3\linewidth}}{
        \cellcolor[rgb]{0.9,0.9,0.9}{
            \makecell[{{p{\linewidth}}}]{
                \texttt{\tiny{[GM$|$GM]}}
                \texttt{COMUNICA: Ho spostato il mio 'P' in (1,1) come concordato. Per quanto riguarda 'S' e 'D', nella mia griglia, 'S' è ancora alla riga 1, colonna 2, e 'D' è alla riga 6, colonna 4. Spostare 'S' alla riga 5 potrebbe essere una buona idea, ma dobbiamo anche considerare la posizione di 'M' e altri oggetti per non creare conflitti. Prima di spostare 'S', vorrei sapere come stanno 'M' e 'A' nella tua griglia adesso, e se ci sono altre mosse che puoi fare per allineare gli oggetti. Inoltre, potresti dirmi se hai già pensato a come spostare 'O' vicino a 'S' e 'D' una volta che li avremo allineati?} \\
            }
        }
    }
    & & \\ \\

    \theutterance \stepcounter{utterance}  
    & & \multicolumn{4}{p{0.6\linewidth}}{
        \cellcolor[rgb]{0.9,0.9,0.9}{
            \makecell[{{p{\linewidth}}}]{
                \texttt{\tiny{[P1$\langle$GM]}}
                \texttt{Il tuo messaggio è stato inoltrato all'altro giocatore.} \\
\\ 
\\ 
\texttt{You are currently playing round 7 of maximum 28.} \\
\texttt{You have collectively accumulated 1 of 16 penalties.} \\
\texttt{L'altro giocatore ha inviato questo messaggio:} \\
\\ 
\\ 
\texttt{"Ho spostato il mio 'P' in (1,1) come concordato. Per quanto riguarda 'S' e 'D', nella mia griglia, 'S' è ancora alla riga 1, colonna 2, e 'D' è alla riga 6, colonna 4. Spostare 'S' alla riga 5 potrebbe essere una buona idea, ma dobbiamo anche considerare la posizione di 'M' e altri oggetti per non creare conflitti. Prima di spostare 'S', vorrei sapere come stanno 'M' e 'A' nella tua griglia adesso, e se ci sono altre mosse che puoi fare per allineare gli oggetti. Inoltre, potresti dirmi se hai già pensato a come spostare 'O' vicino a 'S' e 'D' una volta che li avremo allineati?"} \\
\\ 
\\ 
\texttt{Qual è il tuo prossimo comando?} \\
            }
        }
    }
    & \\ \\

    \theutterance \stepcounter{utterance}  
    & \multicolumn{4}{p{0.6\linewidth}}{
        \cellcolor[rgb]{0.9,0.9,0.9}{
            \makecell[{{p{\linewidth}}}]{
                \texttt{\tiny{[P1$\rangle$GM]}}
                \texttt{COMUNICA: Nella mia griglia, 'M' è ancora alla riga 2, colonna 3, e 'A' è alla riga 3, colonna 1, vicino a 'T', sto pensando di spostare 'S' alla riga 5, colonna 3, per iniziare a creare un cluster di oggetti, e poi potrei spostare 'O' vicino a 'S' una volta che sarà nella nuova posizione, forse alla riga 5, colonna 4, cosa ne pensi di questo piano e come puoi aiutare ad allineare gli oggetti nella tua griglia?} \\
            }
        }
    }
    & & \\ \\

    \theutterance \stepcounter{utterance}  
    & & & \multicolumn{2}{p{0.3\linewidth}}{
        \cellcolor[rgb]{0.9,0.9,0.9}{
            \makecell[{{p{\linewidth}}}]{
                \texttt{\tiny{[GM$|$GM]}}
                \texttt{COMUNICA: Nella mia griglia, 'M' è ancora alla riga 2, colonna 3, e 'A' è alla riga 3, colonna 1, vicino a 'T', sto pensando di spostare 'S' alla riga 5, colonna 3, per iniziare a creare un cluster di oggetti, e poi potrei spostare 'O' vicino a 'S' una volta che sarà nella nuova posizione, forse alla riga 5, colonna 4, cosa ne pensi di questo piano e come puoi aiutare ad allineare gli oggetti nella tua griglia?} \\
            }
        }
    }
    & & \\ \\

    \theutterance \stepcounter{utterance}  
    & & \multicolumn{4}{p{0.6\linewidth}}{
        \cellcolor[rgb]{0.9,0.9,0.9}{
            \makecell[{{p{\linewidth}}}]{
                \texttt{\tiny{[GM$\rangle$P2]}}
                \texttt{Il tuo messaggio è stato inoltrato all'altro giocatore.} \\
\\ 
\\ 
\texttt{You are currently playing round 8 of maximum 28.} \\
\texttt{You have collectively accumulated 1 of 16 penalties.} \\
\texttt{L'altro giocatore ha inviato questo messaggio:} \\
\\ 
\\ 
\texttt{"Nella mia griglia, 'M' è ancora alla riga 2, colonna 3, e 'A' è alla riga 3, colonna 1, vicino a 'T', sto pensando di spostare 'S' alla riga 5, colonna 3, per iniziare a creare un cluster di oggetti, e poi potrei spostare 'O' vicino a 'S' una volta che sarà nella nuova posizione, forse alla riga 5, colonna 4, cosa ne pensi di questo piano e come puoi aiutare ad allineare gli oggetti nella tua griglia?"} \\
\\ 
\\ 
\texttt{Qual è il tuo prossimo comando?} \\
            }
        }
    }
    & \\ \\

    \theutterance \stepcounter{utterance}  
    & & & \multicolumn{4}{p{0.6\linewidth}}{
        \cellcolor[rgb]{0.9,0.9,0.9}{
            \makecell[{{p{\linewidth}}}]{
                \texttt{\tiny{[GM$\langle$P2]}}
                \texttt{COMUNICA: Il tuo piano di spostare 'S' alla riga 5, colonna 3, e poi 'O' alla riga 5, colonna 4, sembra una buona strategia per creare un cluster di oggetti. Per aiutare ad allineare gli oggetti nella mia griglia, potrei spostare 'S' alla riga 5, colonna 3, in modo da allinearla con la tua 'S'. Inoltre, potrei spostare 'D' alla riga 5, colonna 4, in modo da creare un gruppo di oggetti vicini. Prima di fare queste mosse, vorrei confermare che tu sia d'accordo con questo piano e che non ci siano conflitti con gli altri oggetti nella tua griglia. Se sei d'accordo, propongo di eseguire queste mosse nelle prossime tornate.} \\
            }
        }
    }
     \\ \\

    \theutterance \stepcounter{utterance}  
    & & & \multicolumn{2}{p{0.3\linewidth}}{
        \cellcolor[rgb]{0.9,0.9,0.9}{
            \makecell[{{p{\linewidth}}}]{
                \texttt{\tiny{[GM$|$GM]}}
                \texttt{COMUNICA: Il tuo piano di spostare 'S' alla riga 5, colonna 3, e poi 'O' alla riga 5, colonna 4, sembra una buona strategia per creare un cluster di oggetti. Per aiutare ad allineare gli oggetti nella mia griglia, potrei spostare 'S' alla riga 5, colonna 3, in modo da allinearla con la tua 'S'. Inoltre, potrei spostare 'D' alla riga 5, colonna 4, in modo da creare un gruppo di oggetti vicini. Prima di fare queste mosse, vorrei confermare che tu sia d'accordo con questo piano e che non ci siano conflitti con gli altri oggetti nella tua griglia. Se sei d'accordo, propongo di eseguire queste mosse nelle prossime tornate.} \\
            }
        }
    }
    & & \\ \\

    \theutterance \stepcounter{utterance}  
    & & \multicolumn{4}{p{0.6\linewidth}}{
        \cellcolor[rgb]{0.9,0.9,0.9}{
            \makecell[{{p{\linewidth}}}]{
                \texttt{\tiny{[P1$\langle$GM]}}
                \texttt{Il tuo messaggio è stato inoltrato all'altro giocatore.} \\
\\ 
\\ 
\texttt{You are currently playing round 8 of maximum 28.} \\
\texttt{You have collectively accumulated 1 of 16 penalties.} \\
\texttt{L'altro giocatore ha inviato questo messaggio:} \\
\\ 
\\ 
\texttt{"Il tuo piano di spostare 'S' alla riga 5, colonna 3, e poi 'O' alla riga 5, colonna 4, sembra una buona strategia per creare un cluster di oggetti. Per aiutare ad allineare gli oggetti nella mia griglia, potrei spostare 'S' alla riga 5, colonna 3, in modo da allinearla con la tua 'S'. Inoltre, potrei spostare 'D' alla riga 5, colonna 4, in modo da creare un gruppo di oggetti vicini. Prima di fare queste mosse, vorrei confermare che tu sia d'accordo con questo piano e che non ci siano conflitti con gli altri oggetti nella tua griglia. Se sei d'accordo, propongo di eseguire queste mosse nelle prossime tornate."} \\
\\ 
\\ 
\texttt{Qual è il tuo prossimo comando?} \\
            }
        }
    }
    & \\ \\

    \theutterance \stepcounter{utterance}  
    & \multicolumn{4}{p{0.6\linewidth}}{
        \cellcolor[rgb]{0.9,0.9,0.9}{
            \makecell[{{p{\linewidth}}}]{
                \texttt{\tiny{[P1$\rangle$GM]}}
                \texttt{SPOSTA: S, (5, 3)} \\
            }
        }
    }
    & & \\ \\

    \theutterance \stepcounter{utterance}  
    & & & \multicolumn{2}{p{0.3\linewidth}}{
        \cellcolor[rgb]{0.9,0.9,0.9}{
            \makecell[{{p{\linewidth}}}]{
                \texttt{\tiny{[GM$|$GM]}}
                \texttt{SPOSTA: S, (5, 3)} \\
            }
        }
    }
    & & \\ \\

    \theutterance \stepcounter{utterance}  
    & & & \multicolumn{2}{p{0.3\linewidth}}{
        \cellcolor[rgb]{0.9,0.9,0.9}{
            \makecell[{{p{\linewidth}}}]{
                \texttt{\tiny{[GM$|$GM]}}
                \texttt{Penalità: (5,3) non è vuoto, ma contiene '─'.} \\
\\ 
\texttt{You have collectively accumulated 2 of 16 penalties.} \\
\texttt{Riprova!} \\
            }
        }
    }
    & & \\ \\

    \theutterance \stepcounter{utterance}  
    & & \multicolumn{4}{p{0.6\linewidth}}{
        \cellcolor[rgb]{0.9,0.9,0.9}{
            \makecell[{{p{\linewidth}}}]{
                \texttt{\tiny{[P1$\langle$GM]}}
                \texttt{Penalità: (5,3) non è vuoto, ma contiene '─'.} \\
\\ 
\texttt{You have collectively accumulated 2 of 16 penalties.} \\
\texttt{Riprova!} \\
            }
        }
    }
    & \\ \\

    \theutterance \stepcounter{utterance}  
    & \multicolumn{4}{p{0.6\linewidth}}{
        \cellcolor[rgb]{0.9,0.9,0.9}{
            \makecell[{{p{\linewidth}}}]{
                \texttt{\tiny{[P1$\rangle$GM]}}
                \texttt{COMUNICA: Mi scuso per l'errore, ho capito che (5,3) non è vuota, contiene '─', quindi devo ripensare la mossa, forse potrei spostare 'S' in un'altra posizione vicina, come (5,2) o (5,4), ma vorrei discuterne con te prima di decidere, cosa ne pensi?} \\
            }
        }
    }
    & & \\ \\

    \theutterance \stepcounter{utterance}  
    & & & \multicolumn{2}{p{0.3\linewidth}}{
        \cellcolor[rgb]{0.9,0.9,0.9}{
            \makecell[{{p{\linewidth}}}]{
                \texttt{\tiny{[GM$|$GM]}}
                \texttt{COMUNICA: Mi scuso per l'errore, ho capito che (5,3) non è vuota, contiene '─', quindi devo ripensare la mossa, forse potrei spostare 'S' in un'altra posizione vicina, come (5,2) o (5,4), ma vorrei discuterne con te prima di decidere, cosa ne pensi?} \\
            }
        }
    }
    & & \\ \\

    \theutterance \stepcounter{utterance}  
    & & \multicolumn{4}{p{0.6\linewidth}}{
        \cellcolor[rgb]{0.9,0.9,0.9}{
            \makecell[{{p{\linewidth}}}]{
                \texttt{\tiny{[GM$\rangle$P2]}}
                \texttt{Il tuo messaggio è stato inoltrato all'altro giocatore.} \\
\\ 
\\ 
\texttt{You are currently playing round 9 of maximum 28.} \\
\texttt{You have collectively accumulated 2 of 16 penalties.} \\
\texttt{L'altro giocatore ha inviato questo messaggio:} \\
\\ 
\\ 
\texttt{"Mi scuso per l'errore, ho capito che (5,3) non è vuota, contiene '─', quindi devo ripensare la mossa, forse potrei spostare 'S' in un'altra posizione vicina, come (5,2) o (5,4), ma vorrei discuterne con te prima di decidere, cosa ne pensi?"} \\
\\ 
\\ 
\texttt{Qual è il tuo prossimo comando?} \\
            }
        }
    }
    & \\ \\

    \theutterance \stepcounter{utterance}  
    & & & \multicolumn{4}{p{0.6\linewidth}}{
        \cellcolor[rgb]{0.9,0.9,0.9}{
            \makecell[{{p{\linewidth}}}]{
                \texttt{\tiny{[GM$\langle$P2]}}
                \texttt{COMUNICA: Capisco il problema con la posizione (5,3). Spostare 'S' in (5,2) o (5,4) potrebbe essere una buona alternativa. Nella mia griglia, la posizione (5,2) è vuota, quindi potrei spostare il mio 'S' lì se decidessimo di fare questa mossa. La posizione (5,4) è anche una possibilità, ma dovrei controllare se ci sono altri oggetti vicini che potrebbero essere coinvolti. Prima di decidere, vorrei sapere se hai altri oggetti vicini a (5,2) o (5,4) nella tua griglia che potrebbero essere influenzati da questa mossa. Inoltre, potresti dirmi come pensi di gestire la posizione di 'O' una volta che 'S' sarà stato spostato?} \\
            }
        }
    }
     \\ \\

    \theutterance \stepcounter{utterance}  
    & & & \multicolumn{2}{p{0.3\linewidth}}{
        \cellcolor[rgb]{0.9,0.9,0.9}{
            \makecell[{{p{\linewidth}}}]{
                \texttt{\tiny{[GM$|$GM]}}
                \texttt{COMUNICA: Capisco il problema con la posizione (5,3). Spostare 'S' in (5,2) o (5,4) potrebbe essere una buona alternativa. Nella mia griglia, la posizione (5,2) è vuota, quindi potrei spostare il mio 'S' lì se decidessimo di fare questa mossa. La posizione (5,4) è anche una possibilità, ma dovrei controllare se ci sono altri oggetti vicini che potrebbero essere coinvolti. Prima di decidere, vorrei sapere se hai altri oggetti vicini a (5,2) o (5,4) nella tua griglia che potrebbero essere influenzati da questa mossa. Inoltre, potresti dirmi come pensi di gestire la posizione di 'O' una volta che 'S' sarà stato spostato?} \\
            }
        }
    }
    & & \\ \\

    \theutterance \stepcounter{utterance}  
    & & \multicolumn{4}{p{0.6\linewidth}}{
        \cellcolor[rgb]{0.9,0.9,0.9}{
            \makecell[{{p{\linewidth}}}]{
                \texttt{\tiny{[P1$\langle$GM]}}
                \texttt{Il tuo messaggio è stato inoltrato all'altro giocatore.} \\
\\ 
\\ 
\texttt{You are currently playing round 9 of maximum 28.} \\
\texttt{You have collectively accumulated 2 of 16 penalties.} \\
\texttt{L'altro giocatore ha inviato questo messaggio:} \\
\\ 
\\ 
\texttt{"Capisco il problema con la posizione (5,3). Spostare 'S' in (5,2) o (5,4) potrebbe essere una buona alternativa. Nella mia griglia, la posizione (5,2) è vuota, quindi potrei spostare il mio 'S' lì se decidessimo di fare questa mossa. La posizione (5,4) è anche una possibilità, ma dovrei controllare se ci sono altri oggetti vicini che potrebbero essere coinvolti. Prima di decidere, vorrei sapere se hai altri oggetti vicini a (5,2) o (5,4) nella tua griglia che potrebbero essere influenzati da questa mossa. Inoltre, potresti dirmi come pensi di gestire la posizione di 'O' una volta che 'S' sarà stato spostato?"} \\
\\ 
\\ 
\texttt{Qual è il tuo prossimo comando?} \\
            }
        }
    }
    & \\ \\

    \theutterance \stepcounter{utterance}  
    & \multicolumn{4}{p{0.6\linewidth}}{
        \cellcolor[rgb]{0.9,0.9,0.9}{
            \makecell[{{p{\linewidth}}}]{
                \texttt{\tiny{[P1$\rangle$GM]}}
                \texttt{COMUNICA: Nella mia griglia, la posizione (5,2) è vicina a ◌, quindi non ci sono oggetti che potrebbero essere influenzati direttamente, mentre (5,4) è occupata da 'O', quindi spostare 'S' lì potrebbe richiedere ulteriori mosse per riorganizzare 'O', penso che spostare 'S' in (5,2) potrebbe essere la scelta più semplice, e poi potrei spostare 'O' vicino a 'S' una volta che sarà nella nuova posizione, forse in (5,4) o (5,5), cosa ne pensi di questo piano e come puoi aiutare ad allineare gli oggetti nella tua griglia?} \\
            }
        }
    }
    & & \\ \\

    \theutterance \stepcounter{utterance}  
    & & & \multicolumn{2}{p{0.3\linewidth}}{
        \cellcolor[rgb]{0.9,0.9,0.9}{
            \makecell[{{p{\linewidth}}}]{
                \texttt{\tiny{[GM$|$GM]}}
                \texttt{COMUNICA: Nella mia griglia, la posizione (5,2) è vicina a ◌, quindi non ci sono oggetti che potrebbero essere influenzati direttamente, mentre (5,4) è occupata da 'O', quindi spostare 'S' lì potrebbe richiedere ulteriori mosse per riorganizzare 'O', penso che spostare 'S' in (5,2) potrebbe essere la scelta più semplice, e poi potrei spostare 'O' vicino a 'S' una volta che sarà nella nuova posizione, forse in (5,4) o (5,5), cosa ne pensi di questo piano e come puoi aiutare ad allineare gli oggetti nella tua griglia?} \\
            }
        }
    }
    & & \\ \\

    \theutterance \stepcounter{utterance}  
    & & \multicolumn{4}{p{0.6\linewidth}}{
        \cellcolor[rgb]{0.9,0.9,0.9}{
            \makecell[{{p{\linewidth}}}]{
                \texttt{\tiny{[GM$\rangle$P2]}}
                \texttt{Il tuo messaggio è stato inoltrato all'altro giocatore.} \\
\\ 
\\ 
\texttt{You are currently playing round 10 of maximum 28.} \\
\texttt{You have collectively accumulated 2 of 16 penalties.} \\
\texttt{L'altro giocatore ha inviato questo messaggio:} \\
\\ 
\\ 
\texttt{"Nella mia griglia, la posizione (5,2) è vicina a ◌, quindi non ci sono oggetti che potrebbero essere influenzati direttamente, mentre (5,4) è occupata da 'O', quindi spostare 'S' lì potrebbe richiedere ulteriori mosse per riorganizzare 'O', penso che spostare 'S' in (5,2) potrebbe essere la scelta più semplice, e poi potrei spostare 'O' vicino a 'S' una volta che sarà nella nuova posizione, forse in (5,4) o (5,5), cosa ne pensi di questo piano e come puoi aiutare ad allineare gli oggetti nella tua griglia?"} \\
\\ 
\\ 
\texttt{Qual è il tuo prossimo comando?} \\
            }
        }
    }
    & \\ \\

    \theutterance \stepcounter{utterance}  
    & & & \multicolumn{4}{p{0.6\linewidth}}{
        \cellcolor[rgb]{0.9,0.9,0.9}{
            \makecell[{{p{\linewidth}}}]{
                \texttt{\tiny{[GM$\langle$P2]}}
                \texttt{COMUNICA: Il tuo piano di spostare 'S' in (5,2) e poi 'O' vicino a 'S' sembra una buona strategia. Nella mia griglia, potrei spostare il mio 'S' in (5,2) in modo da allinearla con la tua 'S'. Inoltre, potrei spostare il mio 'O' in una posizione che si allinei con la tua 'O' una volta che sarà stata spostata. Prima di eseguire queste mosse, vorrei confermare che tu sia pronto a spostare 'S' in (5,2) e che non ci siano altri oggetti nella tua griglia che potrebbero essere coinvolti in questa mossa. Se sei d'accordo, propongo di eseguire queste mosse nelle prossime tornate. Inoltre, vorrei chiederti se hai già pensato a come allineare 'M', 'A', 'T' e 'D' una volta che 'S' e 'O' saranno nella loro nuova posizione.} \\
            }
        }
    }
     \\ \\

    \theutterance \stepcounter{utterance}  
    & & & \multicolumn{2}{p{0.3\linewidth}}{
        \cellcolor[rgb]{0.9,0.9,0.9}{
            \makecell[{{p{\linewidth}}}]{
                \texttt{\tiny{[GM$|$GM]}}
                \texttt{COMUNICA: Il tuo piano di spostare 'S' in (5,2) e poi 'O' vicino a 'S' sembra una buona strategia. Nella mia griglia, potrei spostare il mio 'S' in (5,2) in modo da allinearla con la tua 'S'. Inoltre, potrei spostare il mio 'O' in una posizione che si allinei con la tua 'O' una volta che sarà stata spostata. Prima di eseguire queste mosse, vorrei confermare che tu sia pronto a spostare 'S' in (5,2) e che non ci siano altri oggetti nella tua griglia che potrebbero essere coinvolti in questa mossa. Se sei d'accordo, propongo di eseguire queste mosse nelle prossime tornate. Inoltre, vorrei chiederti se hai già pensato a come allineare 'M', 'A', 'T' e 'D' una volta che 'S' e 'O' saranno nella loro nuova posizione.} \\
            }
        }
    }
    & & \\ \\

    \theutterance \stepcounter{utterance}  
    & & \multicolumn{4}{p{0.6\linewidth}}{
        \cellcolor[rgb]{0.9,0.9,0.9}{
            \makecell[{{p{\linewidth}}}]{
                \texttt{\tiny{[P1$\langle$GM]}}
                \texttt{Il tuo messaggio è stato inoltrato all'altro giocatore.} \\
\\ 
\\ 
\texttt{You are currently playing round 10 of maximum 28.} \\
\texttt{You have collectively accumulated 2 of 16 penalties.} \\
\texttt{L'altro giocatore ha inviato questo messaggio:} \\
\\ 
\\ 
\texttt{"Il tuo piano di spostare 'S' in (5,2) e poi 'O' vicino a 'S' sembra una buona strategia. Nella mia griglia, potrei spostare il mio 'S' in (5,2) in modo da allinearla con la tua 'S'. Inoltre, potrei spostare il mio 'O' in una posizione che si allinei con la tua 'O' una volta che sarà stata spostata. Prima di eseguire queste mosse, vorrei confermare che tu sia pronto a spostare 'S' in (5,2) e che non ci siano altri oggetti nella tua griglia che potrebbero essere coinvolti in questa mossa. Se sei d'accordo, propongo di eseguire queste mosse nelle prossime tornate. Inoltre, vorrei chiederti se hai già pensato a come allineare 'M', 'A', 'T' e 'D' una volta che 'S' e 'O' saranno nella loro nuova posizione."} \\
\\ 
\\ 
\texttt{Qual è il tuo prossimo comando?} \\
            }
        }
    }
    & \\ \\

    \theutterance \stepcounter{utterance}  
    & \multicolumn{4}{p{0.6\linewidth}}{
        \cellcolor[rgb]{0.9,0.9,0.9}{
            \makecell[{{p{\linewidth}}}]{
                \texttt{\tiny{[P1$\rangle$GM]}}
                \texttt{SPOSTA: S, (5, 2)} \\
            }
        }
    }
    & & \\ \\

    \theutterance \stepcounter{utterance}  
    & & & \multicolumn{2}{p{0.3\linewidth}}{
        \cellcolor[rgb]{0.9,0.9,0.9}{
            \makecell[{{p{\linewidth}}}]{
                \texttt{\tiny{[GM$|$GM]}}
                \texttt{SPOSTA: S, (5, 2)} \\
            }
        }
    }
    & & \\ \\

    \theutterance \stepcounter{utterance}  
    & & & \multicolumn{2}{p{0.3\linewidth}}{
        \cellcolor[rgb]{0.9,0.9,0.9}{
            \makecell[{{p{\linewidth}}}]{
                \texttt{\tiny{[GM$|$GM]}}
                \texttt{Spostato 'S' in (5,2) con successo. La tua griglia aggiornata appare così:} \\
\texttt{```} \\
\texttt{ 1234567} \\
\texttt{╔═══╤═╤═╗} \\
\texttt{║P◌◌│◌│◌║ 1} \\
\texttt{║◌◌M│S├─╢ 2} \\
\texttt{║AT◌├─┘◌║ 3} \\
\texttt{║◌◌◌│◌◌◌║ 4} \\
\texttt{╟─┬─┼───╢ 5} \\
\texttt{║◌│◌│◌◌O║ 6} \\
\texttt{║◌│◌│◌D◌║ 7} \\
\texttt{╚═╧═╧═══╝} \\
\texttt{```} \\
            }
        }
    }
    & & \\ \\

    \theutterance \stepcounter{utterance}  
    & & \multicolumn{4}{p{0.6\linewidth}}{
        \cellcolor[rgb]{0.9,0.9,0.9}{
            \makecell[{{p{\linewidth}}}]{
                \texttt{\tiny{[GM$\rangle$P2]}}
                \texttt{Il tuo messaggio è stato inoltrato all'altro giocatore.} \\
\\ 
\\ 
\texttt{You are currently playing round 11 of maximum 28.} \\
\texttt{You have collectively accumulated 2 of 16 penalties.} \\
\texttt{L'altro giocatore ha spostato un oggetto sulla sua griglia.} \\
\\ 
\\ 
\texttt{Qual è il tuo prossimo comando?} \\
            }
        }
    }
    & \\ \\

    \theutterance \stepcounter{utterance}  
    & & & \multicolumn{4}{p{0.6\linewidth}}{
        \cellcolor[rgb]{0.9,0.9,0.9}{
            \makecell[{{p{\linewidth}}}]{
                \texttt{\tiny{[GM$\langle$P2]}}
                \texttt{SPOSTA: S, (5, 2)} \\
            }
        }
    }
     \\ \\

    \theutterance \stepcounter{utterance}  
    & & & \multicolumn{2}{p{0.3\linewidth}}{
        \cellcolor[rgb]{0.9,0.9,0.9}{
            \makecell[{{p{\linewidth}}}]{
                \texttt{\tiny{[GM$|$GM]}}
                \texttt{SPOSTA: S, (5, 2)} \\
            }
        }
    }
    & & \\ \\

    \theutterance \stepcounter{utterance}  
    & & & \multicolumn{2}{p{0.3\linewidth}}{
        \cellcolor[rgb]{0.9,0.9,0.9}{
            \makecell[{{p{\linewidth}}}]{
                \texttt{\tiny{[GM$|$GM]}}
                \texttt{Spostato 'S' in (5,2) con successo. La tua griglia aggiornata appare così:} \\
\texttt{```} \\
\texttt{ 1234567} \\
\texttt{╔═══╤═╤═╗} \\
\texttt{║PM◌│◌│◌║ 1} \\
\texttt{║◌O◌│S├─╢ 2} \\
\texttt{║◌◌◌├─┘◌║ 3} \\
\texttt{║◌◌◌│◌◌◌║ 4} \\
\texttt{╟─┬─┼───╢ 5} \\
\texttt{║◌│◌│◌◌D║ 6} \\
\texttt{║T│◌│◌◌A║ 7} \\
\texttt{╚═╧═╧═══╝} \\
\texttt{```} \\
            }
        }
    }
    & & \\ \\

    \theutterance \stepcounter{utterance}  
    & & \multicolumn{4}{p{0.6\linewidth}}{
        \cellcolor[rgb]{0.9,0.9,0.9}{
            \makecell[{{p{\linewidth}}}]{
                \texttt{\tiny{[P1$\langle$GM]}}
                \texttt{Spostato 'S' in (5,2) con successo. La tua griglia aggiornata appare così:} \\
\texttt{```} \\
\texttt{ 1234567} \\
\texttt{╔═══╤═╤═╗} \\
\texttt{║P◌◌│◌│◌║ 1} \\
\texttt{║◌◌M│S├─╢ 2} \\
\texttt{║AT◌├─┘◌║ 3} \\
\texttt{║◌◌◌│◌◌◌║ 4} \\
\texttt{╟─┬─┼───╢ 5} \\
\texttt{║◌│◌│◌◌O║ 6} \\
\texttt{║◌│◌│◌D◌║ 7} \\
\texttt{╚═╧═╧═══╝} \\
\texttt{```} \\
\texttt{You are currently playing round 11 of maximum 28.} \\
\texttt{You have collectively accumulated 2 of 16 penalties.} \\
\texttt{L'altro giocatore ha spostato un oggetto sulla sua griglia.} \\
\\ 
\\ 
\texttt{Qual è il tuo prossimo comando?} \\
            }
        }
    }
    & \\ \\

    \theutterance \stepcounter{utterance}  
    & \multicolumn{4}{p{0.6\linewidth}}{
        \cellcolor[rgb]{0.9,0.9,0.9}{
            \makecell[{{p{\linewidth}}}]{
                \texttt{\tiny{[P1$\rangle$GM]}}
                \texttt{SPOSTA: O, (5, 3)} \\
            }
        }
    }
    & & \\ \\

    \theutterance \stepcounter{utterance}  
    & & & \multicolumn{2}{p{0.3\linewidth}}{
        \cellcolor[rgb]{0.9,0.9,0.9}{
            \makecell[{{p{\linewidth}}}]{
                \texttt{\tiny{[GM$|$GM]}}
                \texttt{SPOSTA: O, (5, 3)} \\
            }
        }
    }
    & & \\ \\

    \theutterance \stepcounter{utterance}  
    & & & \multicolumn{2}{p{0.3\linewidth}}{
        \cellcolor[rgb]{0.9,0.9,0.9}{
            \makecell[{{p{\linewidth}}}]{
                \texttt{\tiny{[GM$|$GM]}}
                \texttt{Penalità: (5,3) non è vuoto, ma contiene '─'.} \\
\\ 
\texttt{You have collectively accumulated 3 of 16 penalties.} \\
\texttt{Riprova!} \\
            }
        }
    }
    & & \\ \\

    \theutterance \stepcounter{utterance}  
    & & \multicolumn{4}{p{0.6\linewidth}}{
        \cellcolor[rgb]{0.9,0.9,0.9}{
            \makecell[{{p{\linewidth}}}]{
                \texttt{\tiny{[P1$\langle$GM]}}
                \texttt{Penalità: (5,3) non è vuoto, ma contiene '─'.} \\
\\ 
\texttt{You have collectively accumulated 3 of 16 penalties.} \\
\texttt{Riprova!} \\
            }
        }
    }
    & \\ \\

    \theutterance \stepcounter{utterance}  
    & \multicolumn{4}{p{0.6\linewidth}}{
        \cellcolor[rgb]{0.9,0.9,0.9}{
            \makecell[{{p{\linewidth}}}]{
                \texttt{\tiny{[P1$\rangle$GM]}}
                \texttt{SPOSTA: O, (5, 4)} \\
            }
        }
    }
    & & \\ \\

    \theutterance \stepcounter{utterance}  
    & & & \multicolumn{2}{p{0.3\linewidth}}{
        \cellcolor[rgb]{0.9,0.9,0.9}{
            \makecell[{{p{\linewidth}}}]{
                \texttt{\tiny{[GM$|$GM]}}
                \texttt{SPOSTA: O, (5, 4)} \\
            }
        }
    }
    & & \\ \\

    \theutterance \stepcounter{utterance}  
    & & & \multicolumn{2}{p{0.3\linewidth}}{
        \cellcolor[rgb]{0.9,0.9,0.9}{
            \makecell[{{p{\linewidth}}}]{
                \texttt{\tiny{[GM$|$GM]}}
                \texttt{Spostato 'O' in (5,4) con successo. La tua griglia aggiornata appare così:} \\
\texttt{```} \\
\texttt{ 1234567} \\
\texttt{╔═══╤═╤═╗} \\
\texttt{║P◌◌│◌│◌║ 1} \\
\texttt{║◌◌M│S├─╢ 2} \\
\texttt{║AT◌├─┘◌║ 3} \\
\texttt{║◌◌◌│O◌◌║ 4} \\
\texttt{╟─┬─┼───╢ 5} \\
\texttt{║◌│◌│◌◌◌║ 6} \\
\texttt{║◌│◌│◌D◌║ 7} \\
\texttt{╚═╧═╧═══╝} \\
\texttt{```} \\
            }
        }
    }
    & & \\ \\

    \theutterance \stepcounter{utterance}  
    & & \multicolumn{4}{p{0.6\linewidth}}{
        \cellcolor[rgb]{0.9,0.9,0.9}{
            \makecell[{{p{\linewidth}}}]{
                \texttt{\tiny{[GM$\rangle$P2]}}
                \texttt{Spostato 'S' in (5,2) con successo. La tua griglia aggiornata appare così:} \\
\texttt{```} \\
\texttt{ 1234567} \\
\texttt{╔═══╤═╤═╗} \\
\texttt{║PM◌│◌│◌║ 1} \\
\texttt{║◌O◌│S├─╢ 2} \\
\texttt{║◌◌◌├─┘◌║ 3} \\
\texttt{║◌◌◌│◌◌◌║ 4} \\
\texttt{╟─┬─┼───╢ 5} \\
\texttt{║◌│◌│◌◌D║ 6} \\
\texttt{║T│◌│◌◌A║ 7} \\
\texttt{╚═╧═╧═══╝} \\
\texttt{```} \\
\texttt{You are currently playing round 12 of maximum 28.} \\
\texttt{You have collectively accumulated 3 of 16 penalties.} \\
\texttt{L'altro giocatore ha spostato un oggetto sulla sua griglia.} \\
\\ 
\\ 
\texttt{Qual è il tuo prossimo comando?} \\
            }
        }
    }
    & \\ \\

    \theutterance \stepcounter{utterance}  
    & & & \multicolumn{4}{p{0.6\linewidth}}{
        \cellcolor[rgb]{0.9,0.9,0.9}{
            \makecell[{{p{\linewidth}}}]{
                \texttt{\tiny{[GM$\langle$P2]}}
                \texttt{COMUNICA: Ho spostato il mio 'S' in (5,2) come concordato. Adesso che 'S' è nella sua nuova posizione, vorrei discutere come allineare 'O' e 'D'. Nella mia griglia, 'O' è ancora alla riga 2, colonna 2, e 'D' è alla riga 6, colonna 4. Pensavo di spostare 'D' alla riga 5, colonna 4, per allinearla con 'S', ma vorrei sapere se questo confligge con la tua strategia per 'O' e 'D'. Inoltre, come stanno 'M', 'A' e 'T' nella tua griglia adesso? Ci sono altre mosse che possiamo fare per allineare questi oggetti?} \\
            }
        }
    }
     \\ \\

    \theutterance \stepcounter{utterance}  
    & & & \multicolumn{2}{p{0.3\linewidth}}{
        \cellcolor[rgb]{0.9,0.9,0.9}{
            \makecell[{{p{\linewidth}}}]{
                \texttt{\tiny{[GM$|$GM]}}
                \texttt{COMUNICA: Ho spostato il mio 'S' in (5,2) come concordato. Adesso che 'S' è nella sua nuova posizione, vorrei discutere come allineare 'O' e 'D'. Nella mia griglia, 'O' è ancora alla riga 2, colonna 2, e 'D' è alla riga 6, colonna 4. Pensavo di spostare 'D' alla riga 5, colonna 4, per allinearla con 'S', ma vorrei sapere se questo confligge con la tua strategia per 'O' e 'D'. Inoltre, come stanno 'M', 'A' e 'T' nella tua griglia adesso? Ci sono altre mosse che possiamo fare per allineare questi oggetti?} \\
            }
        }
    }
    & & \\ \\

    \theutterance \stepcounter{utterance}  
    & & \multicolumn{4}{p{0.6\linewidth}}{
        \cellcolor[rgb]{0.9,0.9,0.9}{
            \makecell[{{p{\linewidth}}}]{
                \texttt{\tiny{[P1$\langle$GM]}}
                \texttt{Spostato 'O' in (5,4) con successo. La tua griglia aggiornata appare così:} \\
\texttt{```} \\
\texttt{ 1234567} \\
\texttt{╔═══╤═╤═╗} \\
\texttt{║P◌◌│◌│◌║ 1} \\
\texttt{║◌◌M│S├─╢ 2} \\
\texttt{║AT◌├─┘◌║ 3} \\
\texttt{║◌◌◌│O◌◌║ 4} \\
\texttt{╟─┬─┼───╢ 5} \\
\texttt{║◌│◌│◌◌◌║ 6} \\
\texttt{║◌│◌│◌D◌║ 7} \\
\texttt{╚═╧═╧═══╝} \\
\texttt{```} \\
\texttt{You are currently playing round 12 of maximum 28.} \\
\texttt{You have collectively accumulated 3 of 16 penalties.} \\
\texttt{L'altro giocatore ha inviato questo messaggio:} \\
\\ 
\\ 
\texttt{"Ho spostato il mio 'S' in (5,2) come concordato. Adesso che 'S' è nella sua nuova posizione, vorrei discutere come allineare 'O' e 'D'. Nella mia griglia, 'O' è ancora alla riga 2, colonna 2, e 'D' è alla riga 6, colonna 4. Pensavo di spostare 'D' alla riga 5, colonna 4, per allinearla con 'S', ma vorrei sapere se questo confligge con la tua strategia per 'O' e 'D'. Inoltre, come stanno 'M', 'A' e 'T' nella tua griglia adesso? Ci sono altre mosse che possiamo fare per allineare questi oggetti?"} \\
\\ 
\\ 
\texttt{Qual è il tuo prossimo comando?} \\
            }
        }
    }
    & \\ \\

    \theutterance \stepcounter{utterance}  
    & \multicolumn{4}{p{0.6\linewidth}}{
        \cellcolor[rgb]{0.9,0.9,0.9}{
            \makecell[{{p{\linewidth}}}]{
                \texttt{\tiny{[P1$\rangle$GM]}}
                \texttt{COMUNICA: Nella mia griglia, 'M' è alla riga 2, colonna 3, 'A' è alla riga 3, colonna 1, e 'T' è alla riga 3, colonna 2, quindi potremmo dover fare alcune mosse per allinearli, ma prima vorrei focalizzarmi su 'D', penso che spostare 'D' alla riga 5, colonna 5, potrebbe essere una buona idea, in questo modo potremmo creare un gruppo di oggetti nella parte inferiore della griglia, cosa ne pensi di questo piano e come puoi aiutare ad allineare gli oggetti nella tua griglia?} \\
            }
        }
    }
    & & \\ \\

    \theutterance \stepcounter{utterance}  
    & & & \multicolumn{2}{p{0.3\linewidth}}{
        \cellcolor[rgb]{0.9,0.9,0.9}{
            \makecell[{{p{\linewidth}}}]{
                \texttt{\tiny{[GM$|$GM]}}
                \texttt{COMUNICA: Nella mia griglia, 'M' è alla riga 2, colonna 3, 'A' è alla riga 3, colonna 1, e 'T' è alla riga 3, colonna 2, quindi potremmo dover fare alcune mosse per allinearli, ma prima vorrei focalizzarmi su 'D', penso che spostare 'D' alla riga 5, colonna 5, potrebbe essere una buona idea, in questo modo potremmo creare un gruppo di oggetti nella parte inferiore della griglia, cosa ne pensi di questo piano e come puoi aiutare ad allineare gli oggetti nella tua griglia?} \\
            }
        }
    }
    & & \\ \\

    \theutterance \stepcounter{utterance}  
    & & \multicolumn{4}{p{0.6\linewidth}}{
        \cellcolor[rgb]{0.9,0.9,0.9}{
            \makecell[{{p{\linewidth}}}]{
                \texttt{\tiny{[GM$\rangle$P2]}}
                \texttt{Il tuo messaggio è stato inoltrato all'altro giocatore.} \\
\\ 
\\ 
\texttt{You are currently playing round 13 of maximum 28.} \\
\texttt{You have collectively accumulated 3 of 16 penalties.} \\
\texttt{L'altro giocatore ha inviato questo messaggio:} \\
\\ 
\\ 
\texttt{"Nella mia griglia, 'M' è alla riga 2, colonna 3, 'A' è alla riga 3, colonna 1, e 'T' è alla riga 3, colonna 2, quindi potremmo dover fare alcune mosse per allinearli, ma prima vorrei focalizzarmi su 'D', penso che spostare 'D' alla riga 5, colonna 5, potrebbe essere una buona idea, in questo modo potremmo creare un gruppo di oggetti nella parte inferiore della griglia, cosa ne pensi di questo piano e come puoi aiutare ad allineare gli oggetti nella tua griglia?"} \\
\\ 
\\ 
\texttt{Qual è il tuo prossimo comando?} \\
            }
        }
    }
    & \\ \\

    \theutterance \stepcounter{utterance}  
    & & & \multicolumn{4}{p{0.6\linewidth}}{
        \cellcolor[rgb]{0.9,0.9,0.9}{
            \makecell[{{p{\linewidth}}}]{
                \texttt{\tiny{[GM$\langle$P2]}}
                \texttt{COMUNICA: Il tuo piano di spostare 'D' alla riga 5, colonna 5, sembra una buona strategia per creare un gruppo di oggetti nella parte inferiore della griglia. Nella mia griglia, potrei spostare il mio 'D' alla riga 5, colonna 5, in modo da allinearla con la tua 'D'. Inoltre, potrei valutare la possibilità di spostare 'O' vicino a 'S' e 'D' per creare un cluster di oggetti. Prima di eseguire queste mosse, vorrei confermare che tu sia pronto a spostare 'D' alla riga 5, colonna 5, e che non ci siano altri oggetti nella tua griglia che potrebbero essere coinvolti in questa mossa. Se sei d'accordo, propongo di eseguire queste mosse nelle prossime tornate. Inoltre, vorrei chiederti se hai già pensato a come allineare 'M', 'A' e 'T' una volta che 'D' sarà nella sua nuova posizione.} \\
            }
        }
    }
     \\ \\

    \theutterance \stepcounter{utterance}  
    & & & \multicolumn{2}{p{0.3\linewidth}}{
        \cellcolor[rgb]{0.9,0.9,0.9}{
            \makecell[{{p{\linewidth}}}]{
                \texttt{\tiny{[GM$|$GM]}}
                \texttt{COMUNICA: Il tuo piano di spostare 'D' alla riga 5, colonna 5, sembra una buona strategia per creare un gruppo di oggetti nella parte inferiore della griglia. Nella mia griglia, potrei spostare il mio 'D' alla riga 5, colonna 5, in modo da allinearla con la tua 'D'. Inoltre, potrei valutare la possibilità di spostare 'O' vicino a 'S' e 'D' per creare un cluster di oggetti. Prima di eseguire queste mosse, vorrei confermare che tu sia pronto a spostare 'D' alla riga 5, colonna 5, e che non ci siano altri oggetti nella tua griglia che potrebbero essere coinvolti in questa mossa. Se sei d'accordo, propongo di eseguire queste mosse nelle prossime tornate. Inoltre, vorrei chiederti se hai già pensato a come allineare 'M', 'A' e 'T' una volta che 'D' sarà nella sua nuova posizione.} \\
            }
        }
    }
    & & \\ \\

    \theutterance \stepcounter{utterance}  
    & & \multicolumn{4}{p{0.6\linewidth}}{
        \cellcolor[rgb]{0.9,0.9,0.9}{
            \makecell[{{p{\linewidth}}}]{
                \texttt{\tiny{[P1$\langle$GM]}}
                \texttt{Il tuo messaggio è stato inoltrato all'altro giocatore.} \\
\\ 
\\ 
\texttt{You are currently playing round 13 of maximum 28.} \\
\texttt{You have collectively accumulated 3 of 16 penalties.} \\
\texttt{L'altro giocatore ha inviato questo messaggio:} \\
\\ 
\\ 
\texttt{"Il tuo piano di spostare 'D' alla riga 5, colonna 5, sembra una buona strategia per creare un gruppo di oggetti nella parte inferiore della griglia. Nella mia griglia, potrei spostare il mio 'D' alla riga 5, colonna 5, in modo da allinearla con la tua 'D'. Inoltre, potrei valutare la possibilità di spostare 'O' vicino a 'S' e 'D' per creare un cluster di oggetti. Prima di eseguire queste mosse, vorrei confermare che tu sia pronto a spostare 'D' alla riga 5, colonna 5, e che non ci siano altri oggetti nella tua griglia che potrebbero essere coinvolti in questa mossa. Se sei d'accordo, propongo di eseguire queste mosse nelle prossime tornate. Inoltre, vorrei chiederti se hai già pensato a come allineare 'M', 'A' e 'T' una volta che 'D' sarà nella sua nuova posizione."} \\
\\ 
\\ 
\texttt{Qual è il tuo prossimo comando?} \\
            }
        }
    }
    & \\ \\

    \theutterance \stepcounter{utterance}  
    & \multicolumn{4}{p{0.6\linewidth}}{
        \cellcolor[rgb]{0.9,0.9,0.9}{
            \makecell[{{p{\linewidth}}}]{
                \texttt{\tiny{[P1$\rangle$GM]}}
                \texttt{SPOSTA: D, (5, 5)} \\
            }
        }
    }
    & & \\ \\

    \theutterance \stepcounter{utterance}  
    & & & \multicolumn{2}{p{0.3\linewidth}}{
        \cellcolor[rgb]{0.9,0.9,0.9}{
            \makecell[{{p{\linewidth}}}]{
                \texttt{\tiny{[GM$|$GM]}}
                \texttt{SPOSTA: D, (5, 5)} \\
            }
        }
    }
    & & \\ \\

    \theutterance \stepcounter{utterance}  
    & & & \multicolumn{2}{p{0.3\linewidth}}{
        \cellcolor[rgb]{0.9,0.9,0.9}{
            \makecell[{{p{\linewidth}}}]{
                \texttt{\tiny{[GM$|$GM]}}
                \texttt{Penalità: (5,5) non è vuoto, ma contiene '─'.} \\
\\ 
\texttt{You have collectively accumulated 4 of 16 penalties.} \\
\texttt{Riprova!} \\
            }
        }
    }
    & & \\ \\

    \theutterance \stepcounter{utterance}  
    & & \multicolumn{4}{p{0.6\linewidth}}{
        \cellcolor[rgb]{0.9,0.9,0.9}{
            \makecell[{{p{\linewidth}}}]{
                \texttt{\tiny{[P1$\langle$GM]}}
                \texttt{Penalità: (5,5) non è vuoto, ma contiene '─'.} \\
\\ 
\texttt{You have collectively accumulated 4 of 16 penalties.} \\
\texttt{Riprova!} \\
            }
        }
    }
    & \\ \\

    \theutterance \stepcounter{utterance}  
    & \multicolumn{4}{p{0.6\linewidth}}{
        \cellcolor[rgb]{0.9,0.9,0.9}{
            \makecell[{{p{\linewidth}}}]{
                \texttt{\tiny{[P1$\rangle$GM]}}
                \texttt{SPOSTA: D, (6, 5)} \\
            }
        }
    }
    & & \\ \\

    \theutterance \stepcounter{utterance}  
    & & & \multicolumn{2}{p{0.3\linewidth}}{
        \cellcolor[rgb]{0.9,0.9,0.9}{
            \makecell[{{p{\linewidth}}}]{
                \texttt{\tiny{[GM$|$GM]}}
                \texttt{SPOSTA: D, (6, 5)} \\
            }
        }
    }
    & & \\ \\

    \theutterance \stepcounter{utterance}  
    & & & \multicolumn{2}{p{0.3\linewidth}}{
        \cellcolor[rgb]{0.9,0.9,0.9}{
            \makecell[{{p{\linewidth}}}]{
                \texttt{\tiny{[GM$|$GM]}}
                \texttt{Penalità: (6,5) non è vuoto, ma contiene '─'.} \\
\\ 
\texttt{You have collectively accumulated 5 of 16 penalties.} \\
\texttt{Riprova!} \\
            }
        }
    }
    & & \\ \\

    \theutterance \stepcounter{utterance}  
    & & \multicolumn{4}{p{0.6\linewidth}}{
        \cellcolor[rgb]{0.9,0.9,0.9}{
            \makecell[{{p{\linewidth}}}]{
                \texttt{\tiny{[P1$\langle$GM]}}
                \texttt{Penalità: (6,5) non è vuoto, ma contiene '─'.} \\
\\ 
\texttt{You have collectively accumulated 5 of 16 penalties.} \\
\texttt{Riprova!} \\
            }
        }
    }
    & \\ \\

    \theutterance \stepcounter{utterance}  
    & \multicolumn{4}{p{0.6\linewidth}}{
        \cellcolor[rgb]{0.9,0.9,0.9}{
            \makecell[{{p{\linewidth}}}]{
                \texttt{\tiny{[P1$\rangle$GM]}}
                \texttt{SPOSTA: D, (7, 5)} \\
            }
        }
    }
    & & \\ \\

    \theutterance \stepcounter{utterance}  
    & & & \multicolumn{2}{p{0.3\linewidth}}{
        \cellcolor[rgb]{0.9,0.9,0.9}{
            \makecell[{{p{\linewidth}}}]{
                \texttt{\tiny{[GM$|$GM]}}
                \texttt{SPOSTA: D, (7, 5)} \\
            }
        }
    }
    & & \\ \\

    \theutterance \stepcounter{utterance}  
    & & & \multicolumn{2}{p{0.3\linewidth}}{
        \cellcolor[rgb]{0.9,0.9,0.9}{
            \makecell[{{p{\linewidth}}}]{
                \texttt{\tiny{[GM$|$GM]}}
                \texttt{Penalità: (7,5) non è vuoto, ma contiene '─'.} \\
\\ 
\texttt{You have collectively accumulated 6 of 16 penalties.} \\
\texttt{Riprova!} \\
            }
        }
    }
    & & \\ \\

    \theutterance \stepcounter{utterance}  
    & & \multicolumn{4}{p{0.6\linewidth}}{
        \cellcolor[rgb]{0.9,0.9,0.9}{
            \makecell[{{p{\linewidth}}}]{
                \texttt{\tiny{[P1$\langle$GM]}}
                \texttt{Penalità: (7,5) non è vuoto, ma contiene '─'.} \\
\\ 
\texttt{You have collectively accumulated 6 of 16 penalties.} \\
\texttt{Riprova!} \\
            }
        }
    }
    & \\ \\

    \theutterance \stepcounter{utterance}  
    & \multicolumn{4}{p{0.6\linewidth}}{
        \cellcolor[rgb]{0.9,0.9,0.9}{
            \makecell[{{p{\linewidth}}}]{
                \texttt{\tiny{[P1$\rangle$GM]}}
                \texttt{COMUNICA: Mi scuso per l'errore, sto avendo problemi a trovare una posizione vuota per spostare 'D', forse potremmo rivedere la strategia e trovare un'altra posizione per 'D' che sia vuota e non contenga '─', potresti aiutarmi a trovare una soluzione?} \\
            }
        }
    }
    & & \\ \\

    \theutterance \stepcounter{utterance}  
    & & & \multicolumn{2}{p{0.3\linewidth}}{
        \cellcolor[rgb]{0.9,0.9,0.9}{
            \makecell[{{p{\linewidth}}}]{
                \texttt{\tiny{[GM$|$GM]}}
                \texttt{COMUNICA: Mi scuso per l'errore, sto avendo problemi a trovare una posizione vuota per spostare 'D', forse potremmo rivedere la strategia e trovare un'altra posizione per 'D' che sia vuota e non contenga '─', potresti aiutarmi a trovare una soluzione?} \\
            }
        }
    }
    & & \\ \\

    \theutterance \stepcounter{utterance}  
    & & \multicolumn{4}{p{0.6\linewidth}}{
        \cellcolor[rgb]{0.9,0.9,0.9}{
            \makecell[{{p{\linewidth}}}]{
                \texttt{\tiny{[GM$\rangle$P2]}}
                \texttt{Il tuo messaggio è stato inoltrato all'altro giocatore.} \\
\\ 
\\ 
\texttt{You are currently playing round 14 of maximum 28.} \\
\texttt{You have collectively accumulated 6 of 16 penalties.} \\
\texttt{L'altro giocatore ha inviato questo messaggio:} \\
\\ 
\\ 
\texttt{"Mi scuso per l'errore, sto avendo problemi a trovare una posizione vuota per spostare 'D', forse potremmo rivedere la strategia e trovare un'altra posizione per 'D' che sia vuota e non contenga '─', potresti aiutarmi a trovare una soluzione?"} \\
\\ 
\\ 
\texttt{Qual è il tuo prossimo comando?} \\
            }
        }
    }
    & \\ \\

    \theutterance \stepcounter{utterance}  
    & & & \multicolumn{4}{p{0.6\linewidth}}{
        \cellcolor[rgb]{0.9,0.9,0.9}{
            \makecell[{{p{\linewidth}}}]{
                \texttt{\tiny{[GM$\langle$P2]}}
                \texttt{COMUNICA: Capisco il problema con la posizione di 'D'. Nella mia griglia, potrei valutare le posizioni disponibili per spostare il mio 'D' e trovare una posizione vuota che si allinei con la tua strategia. Una possibilità potrebbe essere spostare 'D' alla riga 5, colonna 4, ma vorrei discuterne con te per assicurarmi che non ci siano conflitti con gli altri oggetti. Inoltre, potrei controllare se ci sono altre posizioni vuote nella mia griglia che potrebbero essere adatte per 'D'. Potresti descrivere le posizioni disponibili nella tua griglia e noi potremmo lavorare insieme per trovare una soluzione?} \\
            }
        }
    }
     \\ \\

    \theutterance \stepcounter{utterance}  
    & & & \multicolumn{2}{p{0.3\linewidth}}{
        \cellcolor[rgb]{0.9,0.9,0.9}{
            \makecell[{{p{\linewidth}}}]{
                \texttt{\tiny{[GM$|$GM]}}
                \texttt{COMUNICA: Capisco il problema con la posizione di 'D'. Nella mia griglia, potrei valutare le posizioni disponibili per spostare il mio 'D' e trovare una posizione vuota che si allinei con la tua strategia. Una possibilità potrebbe essere spostare 'D' alla riga 5, colonna 4, ma vorrei discuterne con te per assicurarmi che non ci siano conflitti con gli altri oggetti. Inoltre, potrei controllare se ci sono altre posizioni vuote nella mia griglia che potrebbero essere adatte per 'D'. Potresti descrivere le posizioni disponibili nella tua griglia e noi potremmo lavorare insieme per trovare una soluzione?} \\
            }
        }
    }
    & & \\ \\

    \theutterance \stepcounter{utterance}  
    & & \multicolumn{4}{p{0.6\linewidth}}{
        \cellcolor[rgb]{0.9,0.9,0.9}{
            \makecell[{{p{\linewidth}}}]{
                \texttt{\tiny{[P1$\langle$GM]}}
                \texttt{Il tuo messaggio è stato inoltrato all'altro giocatore.} \\
\\ 
\\ 
\texttt{You are currently playing round 14 of maximum 28.} \\
\texttt{You have collectively accumulated 6 of 16 penalties.} \\
\texttt{L'altro giocatore ha inviato questo messaggio:} \\
\\ 
\\ 
\texttt{"Capisco il problema con la posizione di 'D'. Nella mia griglia, potrei valutare le posizioni disponibili per spostare il mio 'D' e trovare una posizione vuota che si allinei con la tua strategia. Una possibilità potrebbe essere spostare 'D' alla riga 5, colonna 4, ma vorrei discuterne con te per assicurarmi che non ci siano conflitti con gli altri oggetti. Inoltre, potrei controllare se ci sono altre posizioni vuote nella mia griglia che potrebbero essere adatte per 'D'. Potresti descrivere le posizioni disponibili nella tua griglia e noi potremmo lavorare insieme per trovare una soluzione?"} \\
\\ 
\\ 
\texttt{Qual è il tuo prossimo comando?} \\
            }
        }
    }
    & \\ \\

    \theutterance \stepcounter{utterance}  
    & \multicolumn{4}{p{0.6\linewidth}}{
        \cellcolor[rgb]{0.9,0.9,0.9}{
            \makecell[{{p{\linewidth}}}]{
                \texttt{\tiny{[P1$\rangle$GM]}}
                \texttt{COMUNICA: Nella mia griglia, la posizione (5,4) è già occupata da 'O', ma potrei valutare di spostare 'O' per fare spazio a 'D', oppure potrei considerare altre posizioni come (6,4) o (7,4) per spostare 'D', mi piacerebbe discutere ulteriormente con te per trovare la soluzione migliore, cosa ne pensi di queste opzioni?} \\
            }
        }
    }
    & & \\ \\

    \theutterance \stepcounter{utterance}  
    & & & \multicolumn{2}{p{0.3\linewidth}}{
        \cellcolor[rgb]{0.9,0.9,0.9}{
            \makecell[{{p{\linewidth}}}]{
                \texttt{\tiny{[GM$|$GM]}}
                \texttt{COMUNICA: Nella mia griglia, la posizione (5,4) è già occupata da 'O', ma potrei valutare di spostare 'O' per fare spazio a 'D', oppure potrei considerare altre posizioni come (6,4) o (7,4) per spostare 'D', mi piacerebbe discutere ulteriormente con te per trovare la soluzione migliore, cosa ne pensi di queste opzioni?} \\
            }
        }
    }
    & & \\ \\

    \theutterance \stepcounter{utterance}  
    & & \multicolumn{4}{p{0.6\linewidth}}{
        \cellcolor[rgb]{0.9,0.9,0.9}{
            \makecell[{{p{\linewidth}}}]{
                \texttt{\tiny{[GM$\rangle$P2]}}
                \texttt{Il tuo messaggio è stato inoltrato all'altro giocatore.} \\
\\ 
\\ 
\texttt{You are currently playing round 15 of maximum 28.} \\
\texttt{You have collectively accumulated 6 of 16 penalties.} \\
\texttt{L'altro giocatore ha inviato questo messaggio:} \\
\\ 
\\ 
\texttt{"Nella mia griglia, la posizione (5,4) è già occupata da 'O', ma potrei valutare di spostare 'O' per fare spazio a 'D', oppure potrei considerare altre posizioni come (6,4) o (7,4) per spostare 'D', mi piacerebbe discutere ulteriormente con te per trovare la soluzione migliore, cosa ne pensi di queste opzioni?"} \\
\\ 
\\ 
\texttt{Qual è il tuo prossimo comando?} \\
            }
        }
    }
    & \\ \\

    \theutterance \stepcounter{utterance}  
    & & & \multicolumn{4}{p{0.6\linewidth}}{
        \cellcolor[rgb]{0.9,0.9,0.9}{
            \makecell[{{p{\linewidth}}}]{
                \texttt{\tiny{[GM$\langle$P2]}}
                \texttt{COMUNICA: Le opzioni che hai proposto sembrano interessanti. Nella mia griglia, la posizione (6,4) è attualmente occupata da 'D', quindi potrei valutare di spostarlo in una delle posizioni che hai suggerito. Spostare 'O' per fare spazio a 'D' in (5,4) potrebbe essere una buona opzione, ma vorrei considerare anche le altre possibilità. La posizione (7,4) potrebbe essere una buona scelta, poiché potrebbe aiutare a creare un gruppo di oggetti nella parte inferiore della griglia. Potremmo anche valutare la possibilità di spostare 'D' in (6,5) o (7,5) se queste posizioni sono vuote. Vorrei discutere ulteriormente con te per trovare la soluzione migliore e assicurarmi che non ci siano conflitti con gli altri oggetti. Quale delle opzioni che hai proposto ti sembra la più promettente?} \\
            }
        }
    }
     \\ \\

    \theutterance \stepcounter{utterance}  
    & & & \multicolumn{2}{p{0.3\linewidth}}{
        \cellcolor[rgb]{0.9,0.9,0.9}{
            \makecell[{{p{\linewidth}}}]{
                \texttt{\tiny{[GM$|$GM]}}
                \texttt{COMUNICA: Le opzioni che hai proposto sembrano interessanti. Nella mia griglia, la posizione (6,4) è attualmente occupata da 'D', quindi potrei valutare di spostarlo in una delle posizioni che hai suggerito. Spostare 'O' per fare spazio a 'D' in (5,4) potrebbe essere una buona opzione, ma vorrei considerare anche le altre possibilità. La posizione (7,4) potrebbe essere una buona scelta, poiché potrebbe aiutare a creare un gruppo di oggetti nella parte inferiore della griglia. Potremmo anche valutare la possibilità di spostare 'D' in (6,5) o (7,5) se queste posizioni sono vuote. Vorrei discutere ulteriormente con te per trovare la soluzione migliore e assicurarmi che non ci siano conflitti con gli altri oggetti. Quale delle opzioni che hai proposto ti sembra la più promettente?} \\
            }
        }
    }
    & & \\ \\

    \theutterance \stepcounter{utterance}  
    & & \multicolumn{4}{p{0.6\linewidth}}{
        \cellcolor[rgb]{0.9,0.9,0.9}{
            \makecell[{{p{\linewidth}}}]{
                \texttt{\tiny{[P1$\langle$GM]}}
                \texttt{Il tuo messaggio è stato inoltrato all'altro giocatore.} \\
\\ 
\\ 
\texttt{You are currently playing round 15 of maximum 28.} \\
\texttt{You have collectively accumulated 6 of 16 penalties.} \\
\texttt{L'altro giocatore ha inviato questo messaggio:} \\
\\ 
\\ 
\texttt{"Le opzioni che hai proposto sembrano interessanti. Nella mia griglia, la posizione (6,4) è attualmente occupata da 'D', quindi potrei valutare di spostarlo in una delle posizioni che hai suggerito. Spostare 'O' per fare spazio a 'D' in (5,4) potrebbe essere una buona opzione, ma vorrei considerare anche le altre possibilità. La posizione (7,4) potrebbe essere una buona scelta, poiché potrebbe aiutare a creare un gruppo di oggetti nella parte inferiore della griglia. Potremmo anche valutare la possibilità di spostare 'D' in (6,5) o (7,5) se queste posizioni sono vuote. Vorrei discutere ulteriormente con te per trovare la soluzione migliore e assicurarmi che non ci siano conflitti con gli altri oggetti. Quale delle opzioni che hai proposto ti sembra la più promettente?"} \\
\\ 
\\ 
\texttt{Qual è il tuo prossimo comando?} \\
            }
        }
    }
    & \\ \\

    \theutterance \stepcounter{utterance}  
    & \multicolumn{4}{p{0.6\linewidth}}{
        \cellcolor[rgb]{0.9,0.9,0.9}{
            \makecell[{{p{\linewidth}}}]{
                \texttt{\tiny{[P1$\rangle$GM]}}
                \texttt{COMUNICA: Sono d'accordo che (7,4) potrebbe essere una buona scelta per 'D', poiché aiuterebbe a creare un gruppo di oggetti nella parte inferiore della griglia, e considerando che (7,5) contiene '─', non è una opzione valida, quindi (7,4) sembra la scelta più promettente, potrei spostare 'D' lì e poi valutare come spostare gli altri oggetti per allinearli, cosa ne pensi di procedere con questo piano?} \\
            }
        }
    }
    & & \\ \\

    \theutterance \stepcounter{utterance}  
    & & & \multicolumn{2}{p{0.3\linewidth}}{
        \cellcolor[rgb]{0.9,0.9,0.9}{
            \makecell[{{p{\linewidth}}}]{
                \texttt{\tiny{[GM$|$GM]}}
                \texttt{COMUNICA: Sono d'accordo che (7,4) potrebbe essere una buona scelta per 'D', poiché aiuterebbe a creare un gruppo di oggetti nella parte inferiore della griglia, e considerando che (7,5) contiene '─', non è una opzione valida, quindi (7,4) sembra la scelta più promettente, potrei spostare 'D' lì e poi valutare come spostare gli altri oggetti per allinearli, cosa ne pensi di procedere con questo piano?} \\
            }
        }
    }
    & & \\ \\

    \theutterance \stepcounter{utterance}  
    & & \multicolumn{4}{p{0.6\linewidth}}{
        \cellcolor[rgb]{0.9,0.9,0.9}{
            \makecell[{{p{\linewidth}}}]{
                \texttt{\tiny{[GM$\rangle$P2]}}
                \texttt{Il tuo messaggio è stato inoltrato all'altro giocatore.} \\
\\ 
\\ 
\texttt{You are currently playing round 16 of maximum 28.} \\
\texttt{You have collectively accumulated 6 of 16 penalties.} \\
\texttt{L'altro giocatore ha inviato questo messaggio:} \\
\\ 
\\ 
\texttt{"Sono d'accordo che (7,4) potrebbe essere una buona scelta per 'D', poiché aiuterebbe a creare un gruppo di oggetti nella parte inferiore della griglia, e considerando che (7,5) contiene '─', non è una opzione valida, quindi (7,4) sembra la scelta più promettente, potrei spostare 'D' lì e poi valutare come spostare gli altri oggetti per allinearli, cosa ne pensi di procedere con questo piano?"} \\
\\ 
\\ 
\texttt{Qual è il tuo prossimo comando?} \\
            }
        }
    }
    & \\ \\

    \theutterance \stepcounter{utterance}  
    & & & \multicolumn{4}{p{0.6\linewidth}}{
        \cellcolor[rgb]{0.9,0.9,0.9}{
            \makecell[{{p{\linewidth}}}]{
                \texttt{\tiny{[GM$\langle$P2]}}
                \texttt{SPOSTA: D, (7, 4)} \\
            }
        }
    }
     \\ \\

    \theutterance \stepcounter{utterance}  
    & & & \multicolumn{2}{p{0.3\linewidth}}{
        \cellcolor[rgb]{0.9,0.9,0.9}{
            \makecell[{{p{\linewidth}}}]{
                \texttt{\tiny{[GM$|$GM]}}
                \texttt{SPOSTA: D, (7, 4)} \\
            }
        }
    }
    & & \\ \\

    \theutterance \stepcounter{utterance}  
    & & & \multicolumn{2}{p{0.3\linewidth}}{
        \cellcolor[rgb]{0.9,0.9,0.9}{
            \makecell[{{p{\linewidth}}}]{
                \texttt{\tiny{[GM$|$GM]}}
                \texttt{Spostato 'D' in (7,4) con successo. La tua griglia aggiornata appare così:} \\
\texttt{```} \\
\texttt{ 1234567} \\
\texttt{╔═══╤═╤═╗} \\
\texttt{║PM◌│◌│◌║ 1} \\
\texttt{║◌O◌│S├─╢ 2} \\
\texttt{║◌◌◌├─┘◌║ 3} \\
\texttt{║◌◌◌│◌◌D║ 4} \\
\texttt{╟─┬─┼───╢ 5} \\
\texttt{║◌│◌│◌◌◌║ 6} \\
\texttt{║T│◌│◌◌A║ 7} \\
\texttt{╚═╧═╧═══╝} \\
\texttt{```} \\
            }
        }
    }
    & & \\ \\

    \theutterance \stepcounter{utterance}  
    & & \multicolumn{4}{p{0.6\linewidth}}{
        \cellcolor[rgb]{0.9,0.9,0.9}{
            \makecell[{{p{\linewidth}}}]{
                \texttt{\tiny{[P1$\langle$GM]}}
                \texttt{Il tuo messaggio è stato inoltrato all'altro giocatore.} \\
\\ 
\\ 
\texttt{You are currently playing round 16 of maximum 28.} \\
\texttt{You have collectively accumulated 6 of 16 penalties.} \\
\texttt{L'altro giocatore ha spostato un oggetto sulla sua griglia.} \\
\\ 
\\ 
\texttt{Qual è il tuo prossimo comando?} \\
            }
        }
    }
    & \\ \\

    \theutterance \stepcounter{utterance}  
    & \multicolumn{4}{p{0.6\linewidth}}{
        \cellcolor[rgb]{0.9,0.9,0.9}{
            \makecell[{{p{\linewidth}}}]{
                \texttt{\tiny{[P1$\rangle$GM]}}
                \texttt{SPOSTA: D, (7, 4)} \\
            }
        }
    }
    & & \\ \\

    \theutterance \stepcounter{utterance}  
    & & & \multicolumn{2}{p{0.3\linewidth}}{
        \cellcolor[rgb]{0.9,0.9,0.9}{
            \makecell[{{p{\linewidth}}}]{
                \texttt{\tiny{[GM$|$GM]}}
                \texttt{SPOSTA: D, (7, 4)} \\
            }
        }
    }
    & & \\ \\

    \theutterance \stepcounter{utterance}  
    & & & \multicolumn{2}{p{0.3\linewidth}}{
        \cellcolor[rgb]{0.9,0.9,0.9}{
            \makecell[{{p{\linewidth}}}]{
                \texttt{\tiny{[GM$|$GM]}}
                \texttt{Spostato 'D' in (7,4) con successo. La tua griglia aggiornata appare così:} \\
\texttt{```} \\
\texttt{ 1234567} \\
\texttt{╔═══╤═╤═╗} \\
\texttt{║P◌◌│◌│◌║ 1} \\
\texttt{║◌◌M│S├─╢ 2} \\
\texttt{║AT◌├─┘◌║ 3} \\
\texttt{║◌◌◌│O◌D║ 4} \\
\texttt{╟─┬─┼───╢ 5} \\
\texttt{║◌│◌│◌◌◌║ 6} \\
\texttt{║◌│◌│◌◌◌║ 7} \\
\texttt{╚═╧═╧═══╝} \\
\texttt{```} \\
            }
        }
    }
    & & \\ \\

    \theutterance \stepcounter{utterance}  
    & & \multicolumn{4}{p{0.6\linewidth}}{
        \cellcolor[rgb]{0.9,0.9,0.9}{
            \makecell[{{p{\linewidth}}}]{
                \texttt{\tiny{[GM$\rangle$P2]}}
                \texttt{Spostato 'D' in (7,4) con successo. La tua griglia aggiornata appare così:} \\
\texttt{```} \\
\texttt{ 1234567} \\
\texttt{╔═══╤═╤═╗} \\
\texttt{║PM◌│◌│◌║ 1} \\
\texttt{║◌O◌│S├─╢ 2} \\
\texttt{║◌◌◌├─┘◌║ 3} \\
\texttt{║◌◌◌│◌◌D║ 4} \\
\texttt{╟─┬─┼───╢ 5} \\
\texttt{║◌│◌│◌◌◌║ 6} \\
\texttt{║T│◌│◌◌A║ 7} \\
\texttt{╚═╧═╧═══╝} \\
\texttt{```} \\
\texttt{You are currently playing round 17 of maximum 28.} \\
\texttt{You have collectively accumulated 6 of 16 penalties.} \\
\texttt{L'altro giocatore ha spostato un oggetto sulla sua griglia.} \\
\\ 
\\ 
\texttt{Qual è il tuo prossimo comando?} \\
            }
        }
    }
    & \\ \\

    \theutterance \stepcounter{utterance}  
    & & & \multicolumn{4}{p{0.6\linewidth}}{
        \cellcolor[rgb]{0.9,0.9,0.9}{
            \makecell[{{p{\linewidth}}}]{
                \texttt{\tiny{[GM$\langle$P2]}}
                \texttt{COMUNICA: Ho spostato 'D' in (7,4) come concordato. Adesso che 'D' è nella sua nuova posizione, vorrei discutere come allineare gli altri oggetti. Nella mia griglia, 'O' è ancora alla riga 2, colonna 2, e 'T' è alla riga 7, colonna 1. Pensavo di valutare come spostare 'O' e 'T' per allinearli con gli altri oggetti. Potresti descrivere le posizioni degli oggetti nella tua griglia e come pensi di procedere per allinearli? Inoltre, vorrei chiederti se pensi che siamo vicini a raggiungere l'obiettivo di allineare tutti gli oggetti.} \\
            }
        }
    }
     \\ \\

    \theutterance \stepcounter{utterance}  
    & & & \multicolumn{2}{p{0.3\linewidth}}{
        \cellcolor[rgb]{0.9,0.9,0.9}{
            \makecell[{{p{\linewidth}}}]{
                \texttt{\tiny{[GM$|$GM]}}
                \texttt{COMUNICA: Ho spostato 'D' in (7,4) come concordato. Adesso che 'D' è nella sua nuova posizione, vorrei discutere come allineare gli altri oggetti. Nella mia griglia, 'O' è ancora alla riga 2, colonna 2, e 'T' è alla riga 7, colonna 1. Pensavo di valutare come spostare 'O' e 'T' per allinearli con gli altri oggetti. Potresti descrivere le posizioni degli oggetti nella tua griglia e come pensi di procedere per allinearli? Inoltre, vorrei chiederti se pensi che siamo vicini a raggiungere l'obiettivo di allineare tutti gli oggetti.} \\
            }
        }
    }
    & & \\ \\

    \theutterance \stepcounter{utterance}  
    & & \multicolumn{4}{p{0.6\linewidth}}{
        \cellcolor[rgb]{0.9,0.9,0.9}{
            \makecell[{{p{\linewidth}}}]{
                \texttt{\tiny{[P1$\langle$GM]}}
                \texttt{Spostato 'D' in (7,4) con successo. La tua griglia aggiornata appare così:} \\
\texttt{```} \\
\texttt{ 1234567} \\
\texttt{╔═══╤═╤═╗} \\
\texttt{║P◌◌│◌│◌║ 1} \\
\texttt{║◌◌M│S├─╢ 2} \\
\texttt{║AT◌├─┘◌║ 3} \\
\texttt{║◌◌◌│O◌D║ 4} \\
\texttt{╟─┬─┼───╢ 5} \\
\texttt{║◌│◌│◌◌◌║ 6} \\
\texttt{║◌│◌│◌◌◌║ 7} \\
\texttt{╚═╧═╧═══╝} \\
\texttt{```} \\
\texttt{You are currently playing round 17 of maximum 28.} \\
\texttt{You have collectively accumulated 6 of 16 penalties.} \\
\texttt{L'altro giocatore ha inviato questo messaggio:} \\
\\ 
\\ 
\texttt{"Ho spostato 'D' in (7,4) come concordato. Adesso che 'D' è nella sua nuova posizione, vorrei discutere come allineare gli altri oggetti. Nella mia griglia, 'O' è ancora alla riga 2, colonna 2, e 'T' è alla riga 7, colonna 1. Pensavo di valutare come spostare 'O' e 'T' per allinearli con gli altri oggetti. Potresti descrivere le posizioni degli oggetti nella tua griglia e come pensi di procedere per allinearli? Inoltre, vorrei chiederti se pensi che siamo vicini a raggiungere l'obiettivo di allineare tutti gli oggetti."} \\
\\ 
\\ 
\texttt{Qual è il tuo prossimo comando?} \\
            }
        }
    }
    & \\ \\

    \theutterance \stepcounter{utterance}  
    & \multicolumn{4}{p{0.6\linewidth}}{
        \cellcolor[rgb]{0.9,0.9,0.9}{
            \makecell[{{p{\linewidth}}}]{
                \texttt{\tiny{[P1$\rangle$GM]}}
                \texttt{COMUNICA: Sono felice di vedere che 'D' è stato spostato con successo, nella mia griglia, 'O' è alla riga 4, colonna 4, e 'T' è alla riga 3, colonna 2, quindi forse potremmo valutare di spostare 'T' vicino a 'S' e 'M' per creare un cluster di oggetti, e poi pensare a come spostare 'O' per allinearla con gli altri oggetti, penso che stiamo facendo buoni progressi, ma c'è ancora lavoro da fare, credo che siamo a metà strada per raggiungere l'obiettivo, cosa ne pensi?} \\
            }
        }
    }
    & & \\ \\

    \theutterance \stepcounter{utterance}  
    & & & \multicolumn{2}{p{0.3\linewidth}}{
        \cellcolor[rgb]{0.9,0.9,0.9}{
            \makecell[{{p{\linewidth}}}]{
                \texttt{\tiny{[GM$|$GM]}}
                \texttt{COMUNICA: Sono felice di vedere che 'D' è stato spostato con successo, nella mia griglia, 'O' è alla riga 4, colonna 4, e 'T' è alla riga 3, colonna 2, quindi forse potremmo valutare di spostare 'T' vicino a 'S' e 'M' per creare un cluster di oggetti, e poi pensare a come spostare 'O' per allinearla con gli altri oggetti, penso che stiamo facendo buoni progressi, ma c'è ancora lavoro da fare, credo che siamo a metà strada per raggiungere l'obiettivo, cosa ne pensi?} \\
            }
        }
    }
    & & \\ \\

    \theutterance \stepcounter{utterance}  
    & & \multicolumn{4}{p{0.6\linewidth}}{
        \cellcolor[rgb]{0.9,0.9,0.9}{
            \makecell[{{p{\linewidth}}}]{
                \texttt{\tiny{[GM$\rangle$P2]}}
                \texttt{Il tuo messaggio è stato inoltrato all'altro giocatore.} \\
\\ 
\\ 
\texttt{You are currently playing round 18 of maximum 28.} \\
\texttt{You have collectively accumulated 6 of 16 penalties.} \\
\texttt{L'altro giocatore ha inviato questo messaggio:} \\
\\ 
\\ 
\texttt{"Sono felice di vedere che 'D' è stato spostato con successo, nella mia griglia, 'O' è alla riga 4, colonna 4, e 'T' è alla riga 3, colonna 2, quindi forse potremmo valutare di spostare 'T' vicino a 'S' e 'M' per creare un cluster di oggetti, e poi pensare a come spostare 'O' per allinearla con gli altri oggetti, penso che stiamo facendo buoni progressi, ma c'è ancora lavoro da fare, credo che siamo a metà strada per raggiungere l'obiettivo, cosa ne pensi?"} \\
\\ 
\\ 
\texttt{Qual è il tuo prossimo comando?} \\
            }
        }
    }
    & \\ \\

    \theutterance \stepcounter{utterance}  
    & & & \multicolumn{4}{p{0.6\linewidth}}{
        \cellcolor[rgb]{0.9,0.9,0.9}{
            \makecell[{{p{\linewidth}}}]{
                \texttt{\tiny{[GM$\langle$P2]}}
                \texttt{COMUNICA: Sono d'accordo con te, penso che stiamo facendo buoni progressi e siamo a metà strada per raggiungere l'obiettivo. La tua idea di spostare 'T' vicino a 'S' e 'M' per creare un cluster di oggetti sembra una buona strategia. Nella mia griglia, 'T' è alla riga 7, colonna 1, quindi potrei valutare di spostarlo vicino a 'S' e 'M' per allinearli. Inoltre, potrei pensare a come spostare 'O' per allinearla con gli altri oggetti. Vorrei discutere ulteriormente con te per coordinare le nostre mosse e assicurarmi che non ci siano conflitti tra le nostre strategie. Quale posizione pensi di assegnare a 'T' nella tua griglia per creare un cluster di oggetti? E come pensi di gestire la posizione di 'M' e 'A' in relazione a 'T' e 'S'?} \\
            }
        }
    }
     \\ \\

    \theutterance \stepcounter{utterance}  
    & & & \multicolumn{2}{p{0.3\linewidth}}{
        \cellcolor[rgb]{0.9,0.9,0.9}{
            \makecell[{{p{\linewidth}}}]{
                \texttt{\tiny{[GM$|$GM]}}
                \texttt{COMUNICA: Sono d'accordo con te, penso che stiamo facendo buoni progressi e siamo a metà strada per raggiungere l'obiettivo. La tua idea di spostare 'T' vicino a 'S' e 'M' per creare un cluster di oggetti sembra una buona strategia. Nella mia griglia, 'T' è alla riga 7, colonna 1, quindi potrei valutare di spostarlo vicino a 'S' e 'M' per allinearli. Inoltre, potrei pensare a come spostare 'O' per allinearla con gli altri oggetti. Vorrei discutere ulteriormente con te per coordinare le nostre mosse e assicurarmi che non ci siano conflitti tra le nostre strategie. Quale posizione pensi di assegnare a 'T' nella tua griglia per creare un cluster di oggetti? E come pensi di gestire la posizione di 'M' e 'A' in relazione a 'T' e 'S'?} \\
            }
        }
    }
    & & \\ \\

    \theutterance \stepcounter{utterance}  
    & & \multicolumn{4}{p{0.6\linewidth}}{
        \cellcolor[rgb]{0.9,0.9,0.9}{
            \makecell[{{p{\linewidth}}}]{
                \texttt{\tiny{[P1$\langle$GM]}}
                \texttt{Il tuo messaggio è stato inoltrato all'altro giocatore.} \\
\\ 
\\ 
\texttt{You are currently playing round 18 of maximum 28.} \\
\texttt{You have collectively accumulated 6 of 16 penalties.} \\
\texttt{L'altro giocatore ha inviato questo messaggio:} \\
\\ 
\\ 
\texttt{"Sono d'accordo con te, penso che stiamo facendo buoni progressi e siamo a metà strada per raggiungere l'obiettivo. La tua idea di spostare 'T' vicino a 'S' e 'M' per creare un cluster di oggetti sembra una buona strategia. Nella mia griglia, 'T' è alla riga 7, colonna 1, quindi potrei valutare di spostarlo vicino a 'S' e 'M' per allinearli. Inoltre, potrei pensare a come spostare 'O' per allinearla con gli altri oggetti. Vorrei discutere ulteriormente con te per coordinare le nostre mosse e assicurarmi che non ci siano conflitti tra le nostre strategie. Quale posizione pensi di assegnare a 'T' nella tua griglia per creare un cluster di oggetti? E come pensi di gestire la posizione di 'M' e 'A' in relazione a 'T' e 'S'?"} \\
\\ 
\\ 
\texttt{Qual è il tuo prossimo comando?} \\
            }
        }
    }
    & \\ \\

    \theutterance \stepcounter{utterance}  
    & \multicolumn{4}{p{0.6\linewidth}}{
        \cellcolor[rgb]{0.9,0.9,0.9}{
            \makecell[{{p{\linewidth}}}]{
                \texttt{\tiny{[P1$\rangle$GM]}}
                \texttt{COMUNICA: Penso che assegnare a 'T' la posizione (2,4) potrebbe essere una buona scelta per creare un cluster di oggetti vicino a 'S' e 'M', in questo modo potremmo iniziare a vedere una configurazione più ordinata, per quanto riguarda 'M' e 'A', potrei valutare di spostare 'M' in una posizione vicina a 'T' e 'S', forse (2,5) o (3,5), e 'A' potrebbe essere spostato in una posizione che non interferisca con gli altri oggetti, magari (1,5) o (1,6), ma vorrei discuterne ulteriormente con te per assicurarmi che le nostre strategie siano allineate, cosa ne pensi di queste opzioni?} \\
            }
        }
    }
    & & \\ \\

    \theutterance \stepcounter{utterance}  
    & & & \multicolumn{2}{p{0.3\linewidth}}{
        \cellcolor[rgb]{0.9,0.9,0.9}{
            \makecell[{{p{\linewidth}}}]{
                \texttt{\tiny{[GM$|$GM]}}
                \texttt{COMUNICA: Penso che assegnare a 'T' la posizione (2,4) potrebbe essere una buona scelta per creare un cluster di oggetti vicino a 'S' e 'M', in questo modo potremmo iniziare a vedere una configurazione più ordinata, per quanto riguarda 'M' e 'A', potrei valutare di spostare 'M' in una posizione vicina a 'T' e 'S', forse (2,5) o (3,5), e 'A' potrebbe essere spostato in una posizione che non interferisca con gli altri oggetti, magari (1,5) o (1,6), ma vorrei discuterne ulteriormente con te per assicurarmi che le nostre strategie siano allineate, cosa ne pensi di queste opzioni?} \\
            }
        }
    }
    & & \\ \\

    \theutterance \stepcounter{utterance}  
    & & \multicolumn{4}{p{0.6\linewidth}}{
        \cellcolor[rgb]{0.9,0.9,0.9}{
            \makecell[{{p{\linewidth}}}]{
                \texttt{\tiny{[GM$\rangle$P2]}}
                \texttt{Il tuo messaggio è stato inoltrato all'altro giocatore.} \\
\\ 
\\ 
\texttt{You are currently playing round 19 of maximum 28.} \\
\texttt{You have collectively accumulated 6 of 16 penalties.} \\
\texttt{L'altro giocatore ha inviato questo messaggio:} \\
\\ 
\\ 
\texttt{"Penso che assegnare a 'T' la posizione (2,4) potrebbe essere una buona scelta per creare un cluster di oggetti vicino a 'S' e 'M', in questo modo potremmo iniziare a vedere una configurazione più ordinata, per quanto riguarda 'M' e 'A', potrei valutare di spostare 'M' in una posizione vicina a 'T' e 'S', forse (2,5) o (3,5), e 'A' potrebbe essere spostato in una posizione che non interferisca con gli altri oggetti, magari (1,5) o (1,6), ma vorrei discuterne ulteriormente con te per assicurarmi che le nostre strategie siano allineate, cosa ne pensi di queste opzioni?"} \\
\\ 
\\ 
\texttt{Qual è il tuo prossimo comando?} \\
            }
        }
    }
    & \\ \\

    \theutterance \stepcounter{utterance}  
    & & & \multicolumn{4}{p{0.6\linewidth}}{
        \cellcolor[rgb]{0.9,0.9,0.9}{
            \makecell[{{p{\linewidth}}}]{
                \texttt{\tiny{[GM$\langle$P2]}}
                \texttt{COMUNICA: La tua idea di assegnare a 'T' la posizione (2,4) sembra una buona scelta per creare un cluster di oggetti vicino a 'S' e 'M'. Le opzioni che hai proposto per 'M' e 'A' sembrano ragionevoli, ma vorrei valutare le posizioni disponibili nella mia griglia per assicurarmi che non ci siano conflitti. Nella mia griglia, 'M' è alla riga 1, colonna 1, e 'A' è alla riga 7, colonna 5. Potrei valutare di spostare 'M' in una posizione vicina a 'S' e 'T', ma devo considerare la posizione di 'P' che è già alla riga 1, colonna 1. Per 'A', potrei valutare di spostarlo in una posizione che non interferisca con gli altri oggetti, ma devo considerare la posizione di 'D' che è già alla riga 7, colonna 4. Vorrei discutere ulteriormente con te per trovare una soluzione che funzioni per entrambi. Quale posizione pensi di assegnare a 'O' una volta che 'T', 'M' e 'A' saranno stati spostati?} \\
            }
        }
    }
     \\ \\

    \theutterance \stepcounter{utterance}  
    & & & \multicolumn{2}{p{0.3\linewidth}}{
        \cellcolor[rgb]{0.9,0.9,0.9}{
            \makecell[{{p{\linewidth}}}]{
                \texttt{\tiny{[GM$|$GM]}}
                \texttt{COMUNICA: La tua idea di assegnare a 'T' la posizione (2,4) sembra una buona scelta per creare un cluster di oggetti vicino a 'S' e 'M'. Le opzioni che hai proposto per 'M' e 'A' sembrano ragionevoli, ma vorrei valutare le posizioni disponibili nella mia griglia per assicurarmi che non ci siano conflitti. Nella mia griglia, 'M' è alla riga 1, colonna 1, e 'A' è alla riga 7, colonna 5. Potrei valutare di spostare 'M' in una posizione vicina a 'S' e 'T', ma devo considerare la posizione di 'P' che è già alla riga 1, colonna 1. Per 'A', potrei valutare di spostarlo in una posizione che non interferisca con gli altri oggetti, ma devo considerare la posizione di 'D' che è già alla riga 7, colonna 4. Vorrei discutere ulteriormente con te per trovare una soluzione che funzioni per entrambi. Quale posizione pensi di assegnare a 'O' una volta che 'T', 'M' e 'A' saranno stati spostati?} \\
            }
        }
    }
    & & \\ \\

    \theutterance \stepcounter{utterance}  
    & & \multicolumn{4}{p{0.6\linewidth}}{
        \cellcolor[rgb]{0.9,0.9,0.9}{
            \makecell[{{p{\linewidth}}}]{
                \texttt{\tiny{[P1$\langle$GM]}}
                \texttt{Il tuo messaggio è stato inoltrato all'altro giocatore.} \\
\\ 
\\ 
\texttt{You are currently playing round 19 of maximum 28.} \\
\texttt{You have collectively accumulated 6 of 16 penalties.} \\
\texttt{L'altro giocatore ha inviato questo messaggio:} \\
\\ 
\\ 
\texttt{"La tua idea di assegnare a 'T' la posizione (2,4) sembra una buona scelta per creare un cluster di oggetti vicino a 'S' e 'M'. Le opzioni che hai proposto per 'M' e 'A' sembrano ragionevoli, ma vorrei valutare le posizioni disponibili nella mia griglia per assicurarmi che non ci siano conflitti. Nella mia griglia, 'M' è alla riga 1, colonna 1, e 'A' è alla riga 7, colonna 5. Potrei valutare di spostare 'M' in una posizione vicina a 'S' e 'T', ma devo considerare la posizione di 'P' che è già alla riga 1, colonna 1. Per 'A', potrei valutare di spostarlo in una posizione che non interferisca con gli altri oggetti, ma devo considerare la posizione di 'D' che è già alla riga 7, colonna 4. Vorrei discutere ulteriormente con te per trovare una soluzione che funzioni per entrambi. Quale posizione pensi di assegnare a 'O' una volta che 'T', 'M' e 'A' saranno stati spostati?"} \\
\\ 
\\ 
\texttt{Qual è il tuo prossimo comando?} \\
            }
        }
    }
    & \\ \\

    \theutterance \stepcounter{utterance}  
    & \multicolumn{4}{p{0.6\linewidth}}{
        \cellcolor[rgb]{0.9,0.9,0.9}{
            \makecell[{{p{\linewidth}}}]{
                \texttt{\tiny{[P1$\rangle$GM]}}
                \texttt{COMUNICA: Sono felice di sentire che stiamo lavorando insieme per trovare una soluzione, per 'O', penso che potremmo assegnargli una posizione vicina a 'S' e 'D', forse (6,4) o (7,3), in questo modo potremmo creare un altro cluster di oggetti e iniziare a vedere una configurazione più completa, ma vorrei aspettare di vedere come si sviluppa la situazione con 'T', 'M' e 'A' prima di prendere una decisione definitiva, cosa ne pensi di questo piano e come pensi di gestire la posizione di 'P' nella tua griglia?} \\
            }
        }
    }
    & & \\ \\

    \theutterance \stepcounter{utterance}  
    & & & \multicolumn{2}{p{0.3\linewidth}}{
        \cellcolor[rgb]{0.9,0.9,0.9}{
            \makecell[{{p{\linewidth}}}]{
                \texttt{\tiny{[GM$|$GM]}}
                \texttt{COMUNICA: Sono felice di sentire che stiamo lavorando insieme per trovare una soluzione, per 'O', penso che potremmo assegnargli una posizione vicina a 'S' e 'D', forse (6,4) o (7,3), in questo modo potremmo creare un altro cluster di oggetti e iniziare a vedere una configurazione più completa, ma vorrei aspettare di vedere come si sviluppa la situazione con 'T', 'M' e 'A' prima di prendere una decisione definitiva, cosa ne pensi di questo piano e come pensi di gestire la posizione di 'P' nella tua griglia?} \\
            }
        }
    }
    & & \\ \\

    \theutterance \stepcounter{utterance}  
    & & \multicolumn{4}{p{0.6\linewidth}}{
        \cellcolor[rgb]{0.9,0.9,0.9}{
            \makecell[{{p{\linewidth}}}]{
                \texttt{\tiny{[GM$\rangle$P2]}}
                \texttt{Il tuo messaggio è stato inoltrato all'altro giocatore.} \\
\\ 
\\ 
\texttt{You are currently playing round 20 of maximum 28.} \\
\texttt{You have collectively accumulated 6 of 16 penalties.} \\
\texttt{L'altro giocatore ha inviato questo messaggio:} \\
\\ 
\\ 
\texttt{"Sono felice di sentire che stiamo lavorando insieme per trovare una soluzione, per 'O', penso che potremmo assegnargli una posizione vicina a 'S' e 'D', forse (6,4) o (7,3), in questo modo potremmo creare un altro cluster di oggetti e iniziare a vedere una configurazione più completa, ma vorrei aspettare di vedere come si sviluppa la situazione con 'T', 'M' e 'A' prima di prendere una decisione definitiva, cosa ne pensi di questo piano e come pensi di gestire la posizione di 'P' nella tua griglia?"} \\
\\ 
\\ 
\texttt{Qual è il tuo prossimo comando?} \\
            }
        }
    }
    & \\ \\

    \theutterance \stepcounter{utterance}  
    & & & \multicolumn{4}{p{0.6\linewidth}}{
        \cellcolor[rgb]{0.9,0.9,0.9}{
            \makecell[{{p{\linewidth}}}]{
                \texttt{\tiny{[GM$\langle$P2]}}
                \texttt{COMUNICA: La tua idea di assegnare a 'O' una posizione vicina a 'S' e 'D' sembra una buona strategia per creare un altro cluster di oggetti. Le opzioni che hai proposto per 'O' sembrano ragionevoli, ma vorrei valutare le posizioni disponibili nella mia griglia per assicurarmi che non ci siano conflitti. Per quanto riguarda 'P', nella mia griglia è già alla posizione (1,1), quindi potrebbe essere una buona idea lasciarlo lì come punto di riferimento per gli altri oggetti. Potrei valutare di utilizzare 'P' come ancoraggio per gli altri oggetti e costruire la configurazione intorno a esso. Vorrei discutere ulteriormente con te per trovare una soluzione che funzioni per entrambi. Quale è il tuo piano per gli ultimi round del gioco? Pensi che possiamo raggiungere l'obiettivo di allineare tutti gli oggetti entro i round rimanenti?} \\
            }
        }
    }
     \\ \\

    \theutterance \stepcounter{utterance}  
    & & & \multicolumn{2}{p{0.3\linewidth}}{
        \cellcolor[rgb]{0.9,0.9,0.9}{
            \makecell[{{p{\linewidth}}}]{
                \texttt{\tiny{[GM$|$GM]}}
                \texttt{COMUNICA: La tua idea di assegnare a 'O' una posizione vicina a 'S' e 'D' sembra una buona strategia per creare un altro cluster di oggetti. Le opzioni che hai proposto per 'O' sembrano ragionevoli, ma vorrei valutare le posizioni disponibili nella mia griglia per assicurarmi che non ci siano conflitti. Per quanto riguarda 'P', nella mia griglia è già alla posizione (1,1), quindi potrebbe essere una buona idea lasciarlo lì come punto di riferimento per gli altri oggetti. Potrei valutare di utilizzare 'P' come ancoraggio per gli altri oggetti e costruire la configurazione intorno a esso. Vorrei discutere ulteriormente con te per trovare una soluzione che funzioni per entrambi. Quale è il tuo piano per gli ultimi round del gioco? Pensi che possiamo raggiungere l'obiettivo di allineare tutti gli oggetti entro i round rimanenti?} \\
            }
        }
    }
    & & \\ \\

    \theutterance \stepcounter{utterance}  
    & & \multicolumn{4}{p{0.6\linewidth}}{
        \cellcolor[rgb]{0.9,0.9,0.9}{
            \makecell[{{p{\linewidth}}}]{
                \texttt{\tiny{[P1$\langle$GM]}}
                \texttt{Il tuo messaggio è stato inoltrato all'altro giocatore.} \\
\\ 
\\ 
\texttt{You are currently playing round 20 of maximum 28.} \\
\texttt{You have collectively accumulated 6 of 16 penalties.} \\
\texttt{L'altro giocatore ha inviato questo messaggio:} \\
\\ 
\\ 
\texttt{"La tua idea di assegnare a 'O' una posizione vicina a 'S' e 'D' sembra una buona strategia per creare un altro cluster di oggetti. Le opzioni che hai proposto per 'O' sembrano ragionevoli, ma vorrei valutare le posizioni disponibili nella mia griglia per assicurarmi che non ci siano conflitti. Per quanto riguarda 'P', nella mia griglia è già alla posizione (1,1), quindi potrebbe essere una buona idea lasciarlo lì come punto di riferimento per gli altri oggetti. Potrei valutare di utilizzare 'P' come ancoraggio per gli altri oggetti e costruire la configurazione intorno a esso. Vorrei discutere ulteriormente con te per trovare una soluzione che funzioni per entrambi. Quale è il tuo piano per gli ultimi round del gioco? Pensi che possiamo raggiungere l'obiettivo di allineare tutti gli oggetti entro i round rimanenti?"} \\
\\ 
\\ 
\texttt{Qual è il tuo prossimo comando?} \\
            }
        }
    }
    & \\ \\

    \theutterance \stepcounter{utterance}  
    & \multicolumn{4}{p{0.6\linewidth}}{
        \cellcolor[rgb]{0.9,0.9,0.9}{
            \makecell[{{p{\linewidth}}}]{
                \texttt{\tiny{[P1$\rangle$GM]}}
                \texttt{COMUNICA: Sono d'accordo che utilizzare 'P' come punto di riferimento può essere una buona strategia, e penso che possiamo raggiungere l'obiettivo di allineare tutti gli oggetti entro i round rimanenti, il mio piano è di continuare a lavorare insieme per spostare gli oggetti nelle posizioni che abbiamo discusso, come spostare 'T' vicino a 'S' e 'M', e poi valutare come spostare 'O' per allinearla con gli altri oggetti, credo che con una buona coordinazione e comunicazione possiamo raggiungere l'obiettivo, penso che dovremmo iniziare a pensare a quando chiedere di terminare il gioco, forse quando tutti gli oggetti saranno stati spostati nelle loro posizioni finali, cosa ne pensi di questo piano e sei pronto a iniziare a spostare gli oggetti per raggiungere l'obiettivo?} \\
            }
        }
    }
    & & \\ \\

    \theutterance \stepcounter{utterance}  
    & & & \multicolumn{2}{p{0.3\linewidth}}{
        \cellcolor[rgb]{0.9,0.9,0.9}{
            \makecell[{{p{\linewidth}}}]{
                \texttt{\tiny{[GM$|$GM]}}
                \texttt{COMUNICA: Sono d'accordo che utilizzare 'P' come punto di riferimento può essere una buona strategia, e penso che possiamo raggiungere l'obiettivo di allineare tutti gli oggetti entro i round rimanenti, il mio piano è di continuare a lavorare insieme per spostare gli oggetti nelle posizioni che abbiamo discusso, come spostare 'T' vicino a 'S' e 'M', e poi valutare come spostare 'O' per allinearla con gli altri oggetti, credo che con una buona coordinazione e comunicazione possiamo raggiungere l'obiettivo, penso che dovremmo iniziare a pensare a quando chiedere di terminare il gioco, forse quando tutti gli oggetti saranno stati spostati nelle loro posizioni finali, cosa ne pensi di questo piano e sei pronto a iniziare a spostare gli oggetti per raggiungere l'obiettivo?} \\
            }
        }
    }
    & & \\ \\

    \theutterance \stepcounter{utterance}  
    & & \multicolumn{4}{p{0.6\linewidth}}{
        \cellcolor[rgb]{0.9,0.9,0.9}{
            \makecell[{{p{\linewidth}}}]{
                \texttt{\tiny{[GM$\rangle$P2]}}
                \texttt{Il tuo messaggio è stato inoltrato all'altro giocatore.} \\
\\ 
\\ 
\texttt{You are currently playing round 21 of maximum 28.} \\
\texttt{You have collectively accumulated 6 of 16 penalties.} \\
\texttt{L'altro giocatore ha inviato questo messaggio:} \\
\\ 
\\ 
\texttt{"Sono d'accordo che utilizzare 'P' come punto di riferimento può essere una buona strategia, e penso che possiamo raggiungere l'obiettivo di allineare tutti gli oggetti entro i round rimanenti, il mio piano è di continuare a lavorare insieme per spostare gli oggetti nelle posizioni che abbiamo discusso, come spostare 'T' vicino a 'S' e 'M', e poi valutare come spostare 'O' per allinearla con gli altri oggetti, credo che con una buona coordinazione e comunicazione possiamo raggiungere l'obiettivo, penso che dovremmo iniziare a pensare a quando chiedere di terminare il gioco, forse quando tutti gli oggetti saranno stati spostati nelle loro posizioni finali, cosa ne pensi di questo piano e sei pronto a iniziare a spostare gli oggetti per raggiungere l'obiettivo?"} \\
\\ 
\\ 
\texttt{Qual è il tuo prossimo comando?} \\
            }
        }
    }
    & \\ \\

    \theutterance \stepcounter{utterance}  
    & & & \multicolumn{4}{p{0.6\linewidth}}{
        \cellcolor[rgb]{0.9,0.9,0.9}{
            \makecell[{{p{\linewidth}}}]{
                \texttt{\tiny{[GM$\langle$P2]}}
                \texttt{COMUNICA: Sono d'accordo con il tuo piano e credo che utilizzando 'P' come punto di riferimento e lavorando insieme possiamo raggiungere l'obiettivo di allineare tutti gli oggetti. Sono pronto a iniziare a spostare gli oggetti per raggiungere l'obiettivo. Prima di iniziare, vorrei confermare che siamo d'accordo su quali sono le posizioni finali per ciascun oggetto. Una volta che avremo confermato le posizioni finali, potremo iniziare a spostare gli oggetti e lavorare per raggiungere l'obiettivo. Quando pensi che dovremmo chiedere di terminare il gioco? Dovremmo attendere che tutti gli oggetti siano stati spostati esattamente nelle posizioni finali o possiamo chiedere di terminare il gioco quando saremo vicini all'obiettivo?} \\
            }
        }
    }
     \\ \\

    \theutterance \stepcounter{utterance}  
    & & & \multicolumn{2}{p{0.3\linewidth}}{
        \cellcolor[rgb]{0.9,0.9,0.9}{
            \makecell[{{p{\linewidth}}}]{
                \texttt{\tiny{[GM$|$GM]}}
                \texttt{COMUNICA: Sono d'accordo con il tuo piano e credo che utilizzando 'P' come punto di riferimento e lavorando insieme possiamo raggiungere l'obiettivo di allineare tutti gli oggetti. Sono pronto a iniziare a spostare gli oggetti per raggiungere l'obiettivo. Prima di iniziare, vorrei confermare che siamo d'accordo su quali sono le posizioni finali per ciascun oggetto. Una volta che avremo confermato le posizioni finali, potremo iniziare a spostare gli oggetti e lavorare per raggiungere l'obiettivo. Quando pensi che dovremmo chiedere di terminare il gioco? Dovremmo attendere che tutti gli oggetti siano stati spostati esattamente nelle posizioni finali o possiamo chiedere di terminare il gioco quando saremo vicini all'obiettivo?} \\
            }
        }
    }
    & & \\ \\

    \theutterance \stepcounter{utterance}  
    & & \multicolumn{4}{p{0.6\linewidth}}{
        \cellcolor[rgb]{0.9,0.9,0.9}{
            \makecell[{{p{\linewidth}}}]{
                \texttt{\tiny{[P1$\langle$GM]}}
                \texttt{Il tuo messaggio è stato inoltrato all'altro giocatore.} \\
\\ 
\\ 
\texttt{You are currently playing round 21 of maximum 28.} \\
\texttt{You have collectively accumulated 6 of 16 penalties.} \\
\texttt{L'altro giocatore ha inviato questo messaggio:} \\
\\ 
\\ 
\texttt{"Sono d'accordo con il tuo piano e credo che utilizzando 'P' come punto di riferimento e lavorando insieme possiamo raggiungere l'obiettivo di allineare tutti gli oggetti. Sono pronto a iniziare a spostare gli oggetti per raggiungere l'obiettivo. Prima di iniziare, vorrei confermare che siamo d'accordo su quali sono le posizioni finali per ciascun oggetto. Una volta che avremo confermato le posizioni finali, potremo iniziare a spostare gli oggetti e lavorare per raggiungere l'obiettivo. Quando pensi che dovremmo chiedere di terminare il gioco? Dovremmo attendere che tutti gli oggetti siano stati spostati esattamente nelle posizioni finali o possiamo chiedere di terminare il gioco quando saremo vicini all'obiettivo?"} \\
\\ 
\\ 
\texttt{Qual è il tuo prossimo comando?} \\
            }
        }
    }
    & \\ \\

    \theutterance \stepcounter{utterance}  
    & \multicolumn{4}{p{0.6\linewidth}}{
        \cellcolor[rgb]{0.9,0.9,0.9}{
            \makecell[{{p{\linewidth}}}]{
                \texttt{\tiny{[P1$\rangle$GM]}}
                \texttt{COMUNICA: Sono d'accordo che dovremmo confermare le posizioni finali per ciascun oggetto prima di iniziare a spostarli, penso che dovremmo attendere che tutti gli oggetti siano stati spostati esattamente nelle posizioni finali prima di chiedere di terminare il gioco, in modo da assicurarci di aver raggiunto l'obiettivo con precisione, quando penso che siamo vicini all'obiettivo, potrei inviare un messaggio per chiedere se sei pronto a terminare il gioco, e se si, potremmo entrambi inviare il messaggio "finito" per concludere il gioco, cosa ne pensi di questo piano?} \\
            }
        }
    }
    & & \\ \\

    \theutterance \stepcounter{utterance}  
    & & & \multicolumn{2}{p{0.3\linewidth}}{
        \cellcolor[rgb]{0.9,0.9,0.9}{
            \makecell[{{p{\linewidth}}}]{
                \texttt{\tiny{[GM$|$GM]}}
                \texttt{COMUNICA: Sono d'accordo che dovremmo confermare le posizioni finali per ciascun oggetto prima di iniziare a spostarli, penso che dovremmo attendere che tutti gli oggetti siano stati spostati esattamente nelle posizioni finali prima di chiedere di terminare il gioco, in modo da assicurarci di aver raggiunto l'obiettivo con precisione, quando penso che siamo vicini all'obiettivo, potrei inviare un messaggio per chiedere se sei pronto a terminare il gioco, e se si, potremmo entrambi inviare il messaggio "finito" per concludere il gioco, cosa ne pensi di questo piano?} \\
            }
        }
    }
    & & \\ \\

    \theutterance \stepcounter{utterance}  
    & & \multicolumn{4}{p{0.6\linewidth}}{
        \cellcolor[rgb]{0.9,0.9,0.9}{
            \makecell[{{p{\linewidth}}}]{
                \texttt{\tiny{[GM$\rangle$P2]}}
                \texttt{Il tuo messaggio è stato inoltrato all'altro giocatore.} \\
\\ 
\\ 
\texttt{You are currently playing round 22 of maximum 28.} \\
\texttt{You have collectively accumulated 6 of 16 penalties.} \\
\texttt{L'altro giocatore ha inviato questo messaggio:} \\
\\ 
\\ 
\texttt{"Sono d'accordo che dovremmo confermare le posizioni finali per ciascun oggetto prima di iniziare a spostarli, penso che dovremmo attendere che tutti gli oggetti siano stati spostati esattamente nelle posizioni finali prima di chiedere di terminare il gioco, in modo da assicurarci di aver raggiunto l'obiettivo con precisione, quando penso che siamo vicini all'obiettivo, potrei inviare un messaggio per chiedere se sei pronto a terminare il gioco, e se si, potremmo entrambi inviare il messaggio "finito" per concludere il gioco, cosa ne pensi di questo piano?"} \\
\\ 
\\ 
\texttt{Qual è il tuo prossimo comando?} \\
            }
        }
    }
    & \\ \\

    \theutterance \stepcounter{utterance}  
    & & & \multicolumn{4}{p{0.6\linewidth}}{
        \cellcolor[rgb]{0.9,0.9,0.9}{
            \makecell[{{p{\linewidth}}}]{
                \texttt{\tiny{[GM$\langle$P2]}}
                \texttt{COMUNICA: Sono d'accordo con il tuo piano. Attendere che tutti gli oggetti siano stati spostati esattamente nelle posizioni finali prima di chiedere di terminare il gioco è una buona strategia per assicurarci di aver raggiunto l'obiettivo con precisione. La tua proposta di inviare un messaggio per chiedere se sono pronto a terminare il gioco e poi entrambi inviare il messaggio "finito" per concludere il gioco mi sembra una buona procedura. Sono pronto a procedere con il piano e iniziare a spostare gli oggetti per raggiungere l'obiettivo. Prima di iniziare, vorrei confermare che abbiamo una visione condivisa delle posizioni finali per ciascun oggetto. Potresti fornirmi una panoramica delle posizioni finali che hai in mente per ciascun oggetto?} \\
            }
        }
    }
     \\ \\

    \theutterance \stepcounter{utterance}  
    & & & \multicolumn{2}{p{0.3\linewidth}}{
        \cellcolor[rgb]{0.9,0.9,0.9}{
            \makecell[{{p{\linewidth}}}]{
                \texttt{\tiny{[GM$|$GM]}}
                \texttt{COMUNICA: Sono d'accordo con il tuo piano. Attendere che tutti gli oggetti siano stati spostati esattamente nelle posizioni finali prima di chiedere di terminare il gioco è una buona strategia per assicurarci di aver raggiunto l'obiettivo con precisione. La tua proposta di inviare un messaggio per chiedere se sono pronto a terminare il gioco e poi entrambi inviare il messaggio "finito" per concludere il gioco mi sembra una buona procedura. Sono pronto a procedere con il piano e iniziare a spostare gli oggetti per raggiungere l'obiettivo. Prima di iniziare, vorrei confermare che abbiamo una visione condivisa delle posizioni finali per ciascun oggetto. Potresti fornirmi una panoramica delle posizioni finali che hai in mente per ciascun oggetto?} \\
            }
        }
    }
    & & \\ \\

    \theutterance \stepcounter{utterance}  
    & & \multicolumn{4}{p{0.6\linewidth}}{
        \cellcolor[rgb]{0.9,0.9,0.9}{
            \makecell[{{p{\linewidth}}}]{
                \texttt{\tiny{[P1$\langle$GM]}}
                \texttt{Il tuo messaggio è stato inoltrato all'altro giocatore.} \\
\\ 
\\ 
\texttt{You are currently playing round 22 of maximum 28.} \\
\texttt{You have collectively accumulated 6 of 16 penalties.} \\
\texttt{L'altro giocatore ha inviato questo messaggio:} \\
\\ 
\\ 
\texttt{"Sono d'accordo con il tuo piano. Attendere che tutti gli oggetti siano stati spostati esattamente nelle posizioni finali prima di chiedere di terminare il gioco è una buona strategia per assicurarci di aver raggiunto l'obiettivo con precisione. La tua proposta di inviare un messaggio per chiedere se sono pronto a terminare il gioco e poi entrambi inviare il messaggio "finito" per concludere il gioco mi sembra una buona procedura. Sono pronto a procedere con il piano e iniziare a spostare gli oggetti per raggiungere l'obiettivo. Prima di iniziare, vorrei confermare che abbiamo una visione condivisa delle posizioni finali per ciascun oggetto. Potresti fornirmi una panoramica delle posizioni finali che hai in mente per ciascun oggetto?"} \\
\\ 
\\ 
\texttt{Qual è il tuo prossimo comando?} \\
            }
        }
    }
    & \\ \\

    \theutterance \stepcounter{utterance}  
    & \multicolumn{4}{p{0.6\linewidth}}{
        \cellcolor[rgb]{0.9,0.9,0.9}{
            \makecell[{{p{\linewidth}}}]{
                \texttt{\tiny{[P1$\rangle$GM]}}
                \texttt{COMUNICA: Sono felice di sentire che siamo d'accordo sul piano, per quanto riguarda le posizioni finali, penso che 'P' dovrebbe rimanere in (1,1), 'S' in (5,2), 'D' in (7,4), 'O' potrebbe essere posizionato in (6,4) o (7,3), 'T' potrebbe essere posizionato vicino a 'S' e 'M', forse in (2,4), 'M' potrebbe essere posizionato vicino a 'T' e 'S', forse in (2,5), e 'A' potrebbe essere posizionato in una posizione che non interferisca con gli altri oggetti, forse in (1,5) o (1,6), ma vorrei sentire la tua opinione e vedere se abbiamo una visione condivisa delle posizioni finali, cosa ne pensi di queste posizioni?} \\
            }
        }
    }
    & & \\ \\

    \theutterance \stepcounter{utterance}  
    & & & \multicolumn{2}{p{0.3\linewidth}}{
        \cellcolor[rgb]{0.9,0.9,0.9}{
            \makecell[{{p{\linewidth}}}]{
                \texttt{\tiny{[GM$|$GM]}}
                \texttt{COMUNICA: Sono felice di sentire che siamo d'accordo sul piano, per quanto riguarda le posizioni finali, penso che 'P' dovrebbe rimanere in (1,1), 'S' in (5,2), 'D' in (7,4), 'O' potrebbe essere posizionato in (6,4) o (7,3), 'T' potrebbe essere posizionato vicino a 'S' e 'M', forse in (2,4), 'M' potrebbe essere posizionato vicino a 'T' e 'S', forse in (2,5), e 'A' potrebbe essere posizionato in una posizione che non interferisca con gli altri oggetti, forse in (1,5) o (1,6), ma vorrei sentire la tua opinione e vedere se abbiamo una visione condivisa delle posizioni finali, cosa ne pensi di queste posizioni?} \\
            }
        }
    }
    & & \\ \\

    \theutterance \stepcounter{utterance}  
    & & \multicolumn{4}{p{0.6\linewidth}}{
        \cellcolor[rgb]{0.9,0.9,0.9}{
            \makecell[{{p{\linewidth}}}]{
                \texttt{\tiny{[GM$\rangle$P2]}}
                \texttt{Il tuo messaggio è stato inoltrato all'altro giocatore.} \\
\\ 
\\ 
\texttt{You are currently playing round 23 of maximum 28.} \\
\texttt{You have collectively accumulated 6 of 16 penalties.} \\
\texttt{L'altro giocatore ha inviato questo messaggio:} \\
\\ 
\\ 
\texttt{"Sono felice di sentire che siamo d'accordo sul piano, per quanto riguarda le posizioni finali, penso che 'P' dovrebbe rimanere in (1,1), 'S' in (5,2), 'D' in (7,4), 'O' potrebbe essere posizionato in (6,4) o (7,3), 'T' potrebbe essere posizionato vicino a 'S' e 'M', forse in (2,4), 'M' potrebbe essere posizionato vicino a 'T' e 'S', forse in (2,5), e 'A' potrebbe essere posizionato in una posizione che non interferisca con gli altri oggetti, forse in (1,5) o (1,6), ma vorrei sentire la tua opinione e vedere se abbiamo una visione condivisa delle posizioni finali, cosa ne pensi di queste posizioni?"} \\
\\ 
\\ 
\texttt{Qual è il tuo prossimo comando?} \\
            }
        }
    }
    & \\ \\

    \theutterance \stepcounter{utterance}  
    & & & \multicolumn{4}{p{0.6\linewidth}}{
        \cellcolor[rgb]{0.9,0.9,0.9}{
            \makecell[{{p{\linewidth}}}]{
                \texttt{\tiny{[GM$\langle$P2]}}
                \texttt{COMUNICA: Sono d'accordo con le posizioni finali che hai proposto per gli oggetti. Penso che 'P' in (1,1) sia una buona scelta come punto di riferimento, e 'S' in (5,2) sembra una buona posizione per creare un cluster di oggetti. 'D' in (7,4) è già nella posizione corretta. Per 'O', penso che (6,4) o (7,3) siano entrambe buone opzioni, ma potremmo valutare quale sia la migliore in base alle mosse successive. La posizione di 'T' vicino a 'S' e 'M' in (2,4) sembra una buona scelta, e 'M' vicino a 'T' e 'S' in (2,5) è una buona opzione. Per 'A', penso che (1,5) o (1,6) siano entrambe buone opzioni, ma potremmo valutare quale sia la migliore in base alle mosse successive. In generale, penso che abbiamo una visione condivisa delle posizioni finali e possiamo iniziare a lavorare per raggiungere l'obiettivo. Prima di iniziare a spostare gli oggetti, vorrei confermare che siamo d'accordo su quali oggetti spostare per primi e in che ordine. Quale oggetto pensi di spostare per primo?} \\
            }
        }
    }
     \\ \\

    \theutterance \stepcounter{utterance}  
    & & & \multicolumn{2}{p{0.3\linewidth}}{
        \cellcolor[rgb]{0.9,0.9,0.9}{
            \makecell[{{p{\linewidth}}}]{
                \texttt{\tiny{[GM$|$GM]}}
                \texttt{COMUNICA: Sono d'accordo con le posizioni finali che hai proposto per gli oggetti. Penso che 'P' in (1,1) sia una buona scelta come punto di riferimento, e 'S' in (5,2) sembra una buona posizione per creare un cluster di oggetti. 'D' in (7,4) è già nella posizione corretta. Per 'O', penso che (6,4) o (7,3) siano entrambe buone opzioni, ma potremmo valutare quale sia la migliore in base alle mosse successive. La posizione di 'T' vicino a 'S' e 'M' in (2,4) sembra una buona scelta, e 'M' vicino a 'T' e 'S' in (2,5) è una buona opzione. Per 'A', penso che (1,5) o (1,6) siano entrambe buone opzioni, ma potremmo valutare quale sia la migliore in base alle mosse successive. In generale, penso che abbiamo una visione condivisa delle posizioni finali e possiamo iniziare a lavorare per raggiungere l'obiettivo. Prima di iniziare a spostare gli oggetti, vorrei confermare che siamo d'accordo su quali oggetti spostare per primi e in che ordine. Quale oggetto pensi di spostare per primo?} \\
            }
        }
    }
    & & \\ \\

    \theutterance \stepcounter{utterance}  
    & & \multicolumn{4}{p{0.6\linewidth}}{
        \cellcolor[rgb]{0.9,0.9,0.9}{
            \makecell[{{p{\linewidth}}}]{
                \texttt{\tiny{[P1$\langle$GM]}}
                \texttt{Il tuo messaggio è stato inoltrato all'altro giocatore.} \\
\\ 
\\ 
\texttt{You are currently playing round 23 of maximum 28.} \\
\texttt{You have collectively accumulated 6 of 16 penalties.} \\
\texttt{L'altro giocatore ha inviato questo messaggio:} \\
\\ 
\\ 
\texttt{"Sono d'accordo con le posizioni finali che hai proposto per gli oggetti. Penso che 'P' in (1,1) sia una buona scelta come punto di riferimento, e 'S' in (5,2) sembra una buona posizione per creare un cluster di oggetti. 'D' in (7,4) è già nella posizione corretta. Per 'O', penso che (6,4) o (7,3) siano entrambe buone opzioni, ma potremmo valutare quale sia la migliore in base alle mosse successive. La posizione di 'T' vicino a 'S' e 'M' in (2,4) sembra una buona scelta, e 'M' vicino a 'T' e 'S' in (2,5) è una buona opzione. Per 'A', penso che (1,5) o (1,6) siano entrambe buone opzioni, ma potremmo valutare quale sia la migliore in base alle mosse successive. In generale, penso che abbiamo una visione condivisa delle posizioni finali e possiamo iniziare a lavorare per raggiungere l'obiettivo. Prima di iniziare a spostare gli oggetti, vorrei confermare che siamo d'accordo su quali oggetti spostare per primi e in che ordine. Quale oggetto pensi di spostare per primo?"} \\
\\ 
\\ 
\texttt{Qual è il tuo prossimo comando?} \\
            }
        }
    }
    & \\ \\

    \theutterance \stepcounter{utterance}  
    & \multicolumn{4}{p{0.6\linewidth}}{
        \cellcolor[rgb]{0.9,0.9,0.9}{
            \makecell[{{p{\linewidth}}}]{
                \texttt{\tiny{[P1$\rangle$GM]}}
                \texttt{COMUNICA: Sono felice di sentire che siamo d'accordo sulle posizioni finali, per quanto riguarda l'ordine degli spostamenti, penso che potremmo iniziare spostando 'T' in (2,4) per creare un cluster di oggetti con 'S' e 'M', poi potremmo valutare come spostare 'O' in (6,4) o (7,3) e 'A' in (1,5) o (1,6), in base alle esigenze della griglia, credo che spostare 'T' per primo ci darà una buona base per le mosse successive, cosa ne pensi di iniziare con questo spostamento?} \\
            }
        }
    }
    & & \\ \\

    \theutterance \stepcounter{utterance}  
    & & & \multicolumn{2}{p{0.3\linewidth}}{
        \cellcolor[rgb]{0.9,0.9,0.9}{
            \makecell[{{p{\linewidth}}}]{
                \texttt{\tiny{[GM$|$GM]}}
                \texttt{COMUNICA: Sono felice di sentire che siamo d'accordo sulle posizioni finali, per quanto riguarda l'ordine degli spostamenti, penso che potremmo iniziare spostando 'T' in (2,4) per creare un cluster di oggetti con 'S' e 'M', poi potremmo valutare come spostare 'O' in (6,4) o (7,3) e 'A' in (1,5) o (1,6), in base alle esigenze della griglia, credo che spostare 'T' per primo ci darà una buona base per le mosse successive, cosa ne pensi di iniziare con questo spostamento?} \\
            }
        }
    }
    & & \\ \\

    \theutterance \stepcounter{utterance}  
    & & \multicolumn{4}{p{0.6\linewidth}}{
        \cellcolor[rgb]{0.9,0.9,0.9}{
            \makecell[{{p{\linewidth}}}]{
                \texttt{\tiny{[GM$\rangle$P2]}}
                \texttt{Il tuo messaggio è stato inoltrato all'altro giocatore.} \\
\\ 
\\ 
\texttt{You are currently playing round 24 of maximum 28.} \\
\texttt{You have collectively accumulated 6 of 16 penalties.} \\
\texttt{L'altro giocatore ha inviato questo messaggio:} \\
\\ 
\\ 
\texttt{"Sono felice di sentire che siamo d'accordo sulle posizioni finali, per quanto riguarda l'ordine degli spostamenti, penso che potremmo iniziare spostando 'T' in (2,4) per creare un cluster di oggetti con 'S' e 'M', poi potremmo valutare come spostare 'O' in (6,4) o (7,3) e 'A' in (1,5) o (1,6), in base alle esigenze della griglia, credo che spostare 'T' per primo ci darà una buona base per le mosse successive, cosa ne pensi di iniziare con questo spostamento?"} \\
\\ 
\\ 
\texttt{Qual è il tuo prossimo comando?} \\
            }
        }
    }
    & \\ \\

    \theutterance \stepcounter{utterance}  
    & & & \multicolumn{4}{p{0.6\linewidth}}{
        \cellcolor[rgb]{0.9,0.9,0.9}{
            \makecell[{{p{\linewidth}}}]{
                \texttt{\tiny{[GM$\langle$P2]}}
                \texttt{SPOSTA: T, (2, 4)} \\
            }
        }
    }
     \\ \\

    \theutterance \stepcounter{utterance}  
    & & & \multicolumn{2}{p{0.3\linewidth}}{
        \cellcolor[rgb]{0.9,0.9,0.9}{
            \makecell[{{p{\linewidth}}}]{
                \texttt{\tiny{[GM$|$GM]}}
                \texttt{SPOSTA: T, (2, 4)} \\
            }
        }
    }
    & & \\ \\

    \theutterance \stepcounter{utterance}  
    & & & \multicolumn{2}{p{0.3\linewidth}}{
        \cellcolor[rgb]{0.9,0.9,0.9}{
            \makecell[{{p{\linewidth}}}]{
                \texttt{\tiny{[GM$|$GM]}}
                \texttt{Spostato 'T' in (2,4) con successo. La tua griglia aggiornata appare così:} \\
\texttt{```} \\
\texttt{ 1234567} \\
\texttt{╔═══╤═╤═╗} \\
\texttt{║PM◌│◌│◌║ 1} \\
\texttt{║◌O◌│S├─╢ 2} \\
\texttt{║◌◌◌├─┘◌║ 3} \\
\texttt{║◌T◌│◌◌D║ 4} \\
\texttt{╟─┬─┼───╢ 5} \\
\texttt{║◌│◌│◌◌◌║ 6} \\
\texttt{║◌│◌│◌◌A║ 7} \\
\texttt{╚═╧═╧═══╝} \\
\texttt{```} \\
            }
        }
    }
    & & \\ \\

    \theutterance \stepcounter{utterance}  
    & & \multicolumn{4}{p{0.6\linewidth}}{
        \cellcolor[rgb]{0.9,0.9,0.9}{
            \makecell[{{p{\linewidth}}}]{
                \texttt{\tiny{[P1$\langle$GM]}}
                \texttt{Il tuo messaggio è stato inoltrato all'altro giocatore.} \\
\\ 
\\ 
\texttt{You are currently playing round 24 of maximum 28.} \\
\texttt{You have collectively accumulated 6 of 16 penalties.} \\
\texttt{L'altro giocatore ha spostato un oggetto sulla sua griglia.} \\
\\ 
\\ 
\texttt{Qual è il tuo prossimo comando?} \\
            }
        }
    }
    & \\ \\

    \theutterance \stepcounter{utterance}  
    & \multicolumn{4}{p{0.6\linewidth}}{
        \cellcolor[rgb]{0.9,0.9,0.9}{
            \makecell[{{p{\linewidth}}}]{
                \texttt{\tiny{[P1$\rangle$GM]}}
                \texttt{SPOSTA: T, (2, 4)} \\
            }
        }
    }
    & & \\ \\

    \theutterance \stepcounter{utterance}  
    & & & \multicolumn{2}{p{0.3\linewidth}}{
        \cellcolor[rgb]{0.9,0.9,0.9}{
            \makecell[{{p{\linewidth}}}]{
                \texttt{\tiny{[GM$|$GM]}}
                \texttt{SPOSTA: T, (2, 4)} \\
            }
        }
    }
    & & \\ \\

    \theutterance \stepcounter{utterance}  
    & & & \multicolumn{2}{p{0.3\linewidth}}{
        \cellcolor[rgb]{0.9,0.9,0.9}{
            \makecell[{{p{\linewidth}}}]{
                \texttt{\tiny{[GM$|$GM]}}
                \texttt{Spostato 'T' in (2,4) con successo. La tua griglia aggiornata appare così:} \\
\texttt{```} \\
\texttt{ 1234567} \\
\texttt{╔═══╤═╤═╗} \\
\texttt{║P◌◌│◌│◌║ 1} \\
\texttt{║◌◌M│S├─╢ 2} \\
\texttt{║A◌◌├─┘◌║ 3} \\
\texttt{║◌T◌│O◌D║ 4} \\
\texttt{╟─┬─┼───╢ 5} \\
\texttt{║◌│◌│◌◌◌║ 6} \\
\texttt{║◌│◌│◌◌◌║ 7} \\
\texttt{╚═╧═╧═══╝} \\
\texttt{```} \\
            }
        }
    }
    & & \\ \\

    \theutterance \stepcounter{utterance}  
    & & \multicolumn{4}{p{0.6\linewidth}}{
        \cellcolor[rgb]{0.9,0.9,0.9}{
            \makecell[{{p{\linewidth}}}]{
                \texttt{\tiny{[GM$\rangle$P2]}}
                \texttt{Spostato 'T' in (2,4) con successo. La tua griglia aggiornata appare così:} \\
\texttt{```} \\
\texttt{ 1234567} \\
\texttt{╔═══╤═╤═╗} \\
\texttt{║PM◌│◌│◌║ 1} \\
\texttt{║◌O◌│S├─╢ 2} \\
\texttt{║◌◌◌├─┘◌║ 3} \\
\texttt{║◌T◌│◌◌D║ 4} \\
\texttt{╟─┬─┼───╢ 5} \\
\texttt{║◌│◌│◌◌◌║ 6} \\
\texttt{║◌│◌│◌◌A║ 7} \\
\texttt{╚═╧═╧═══╝} \\
\texttt{```} \\
\texttt{You are currently playing round 25 of maximum 28.} \\
\texttt{You have collectively accumulated 6 of 16 penalties.} \\
\texttt{L'altro giocatore ha spostato un oggetto sulla sua griglia.} \\
\\ 
\\ 
\texttt{Qual è il tuo prossimo comando?} \\
            }
        }
    }
    & \\ \\

    \theutterance \stepcounter{utterance}  
    & & & \multicolumn{4}{p{0.6\linewidth}}{
        \cellcolor[rgb]{0.9,0.9,0.9}{
            \makecell[{{p{\linewidth}}}]{
                \texttt{\tiny{[GM$\langle$P2]}}
                \texttt{COMUNICA: Ho spostato 'T' in (2,4) come concordato. Adesso che 'T' è nella sua nuova posizione, vorrei valutare come spostare 'O' e 'A' per allinearli con gli altri oggetti. Nella mia griglia, 'O' è ancora alla riga 2, colonna 2, e 'A' è alla riga 7, colonna 5. Penso che potremmo spostare 'O' in (6,4) o (7,3) e 'A' in (1,5) o (1,6), come discusso in precedenza. Vorrei discutere ulteriormente con te per determinare le migliori mosse successive. Come stanno 'O' e 'A' nella tua griglia e quali sono le tue proposte per spostarli?} \\
            }
        }
    }
     \\ \\

    \theutterance \stepcounter{utterance}  
    & & & \multicolumn{2}{p{0.3\linewidth}}{
        \cellcolor[rgb]{0.9,0.9,0.9}{
            \makecell[{{p{\linewidth}}}]{
                \texttt{\tiny{[GM$|$GM]}}
                \texttt{COMUNICA: Ho spostato 'T' in (2,4) come concordato. Adesso che 'T' è nella sua nuova posizione, vorrei valutare come spostare 'O' e 'A' per allinearli con gli altri oggetti. Nella mia griglia, 'O' è ancora alla riga 2, colonna 2, e 'A' è alla riga 7, colonna 5. Penso che potremmo spostare 'O' in (6,4) o (7,3) e 'A' in (1,5) o (1,6), come discusso in precedenza. Vorrei discutere ulteriormente con te per determinare le migliori mosse successive. Come stanno 'O' e 'A' nella tua griglia e quali sono le tue proposte per spostarli?} \\
            }
        }
    }
    & & \\ \\

    \theutterance \stepcounter{utterance}  
    & & \multicolumn{4}{p{0.6\linewidth}}{
        \cellcolor[rgb]{0.9,0.9,0.9}{
            \makecell[{{p{\linewidth}}}]{
                \texttt{\tiny{[P1$\langle$GM]}}
                \texttt{Spostato 'T' in (2,4) con successo. La tua griglia aggiornata appare così:} \\
\texttt{```} \\
\texttt{ 1234567} \\
\texttt{╔═══╤═╤═╗} \\
\texttt{║P◌◌│◌│◌║ 1} \\
\texttt{║◌◌M│S├─╢ 2} \\
\texttt{║A◌◌├─┘◌║ 3} \\
\texttt{║◌T◌│O◌D║ 4} \\
\texttt{╟─┬─┼───╢ 5} \\
\texttt{║◌│◌│◌◌◌║ 6} \\
\texttt{║◌│◌│◌◌◌║ 7} \\
\texttt{╚═╧═╧═══╝} \\
\texttt{```} \\
\texttt{You are currently playing round 25 of maximum 28.} \\
\texttt{You have collectively accumulated 6 of 16 penalties.} \\
\texttt{L'altro giocatore ha inviato questo messaggio:} \\
\\ 
\\ 
\texttt{"Ho spostato 'T' in (2,4) come concordato. Adesso che 'T' è nella sua nuova posizione, vorrei valutare come spostare 'O' e 'A' per allinearli con gli altri oggetti. Nella mia griglia, 'O' è ancora alla riga 2, colonna 2, e 'A' è alla riga 7, colonna 5. Penso che potremmo spostare 'O' in (6,4) o (7,3) e 'A' in (1,5) o (1,6), come discusso in precedenza. Vorrei discutere ulteriormente con te per determinare le migliori mosse successive. Come stanno 'O' e 'A' nella tua griglia e quali sono le tue proposte per spostarli?"} \\
\\ 
\\ 
\texttt{Qual è il tuo prossimo comando?} \\
            }
        }
    }
    & \\ \\

    \theutterance \stepcounter{utterance}  
    & \multicolumn{4}{p{0.6\linewidth}}{
        \cellcolor[rgb]{0.9,0.9,0.9}{
            \makecell[{{p{\linewidth}}}]{
                \texttt{\tiny{[P1$\rangle$GM]}}
                \texttt{COMUNICA: Nella mia griglia, 'O' è attualmente in (4,4) e 'A' è in (3,1), quindi penso che potremmo spostare 'O' in (6,4) e 'A' in (1,5) per allinearli con gli altri oggetti, ma vorrei discuterne ulteriormente con te per assicurarmi che non ci siano conflitti con le tue mosse, cosa ne pensi di queste proposte e come pensi di gestire le posizioni di 'O' e 'A' nella tua griglia?} \\
            }
        }
    }
    & & \\ \\

    \theutterance \stepcounter{utterance}  
    & & & \multicolumn{2}{p{0.3\linewidth}}{
        \cellcolor[rgb]{0.9,0.9,0.9}{
            \makecell[{{p{\linewidth}}}]{
                \texttt{\tiny{[GM$|$GM]}}
                \texttt{COMUNICA: Nella mia griglia, 'O' è attualmente in (4,4) e 'A' è in (3,1), quindi penso che potremmo spostare 'O' in (6,4) e 'A' in (1,5) per allinearli con gli altri oggetti, ma vorrei discuterne ulteriormente con te per assicurarmi che non ci siano conflitti con le tue mosse, cosa ne pensi di queste proposte e come pensi di gestire le posizioni di 'O' e 'A' nella tua griglia?} \\
            }
        }
    }
    & & \\ \\

    \theutterance \stepcounter{utterance}  
    & & \multicolumn{4}{p{0.6\linewidth}}{
        \cellcolor[rgb]{0.9,0.9,0.9}{
            \makecell[{{p{\linewidth}}}]{
                \texttt{\tiny{[GM$\rangle$P2]}}
                \texttt{Il tuo messaggio è stato inoltrato all'altro giocatore.} \\
\\ 
\\ 
\texttt{You are currently playing round 26 of maximum 28.} \\
\texttt{You have collectively accumulated 6 of 16 penalties.} \\
\texttt{L'altro giocatore ha inviato questo messaggio:} \\
\\ 
\\ 
\texttt{"Nella mia griglia, 'O' è attualmente in (4,4) e 'A' è in (3,1), quindi penso che potremmo spostare 'O' in (6,4) e 'A' in (1,5) per allinearli con gli altri oggetti, ma vorrei discuterne ulteriormente con te per assicurarmi che non ci siano conflitti con le tue mosse, cosa ne pensi di queste proposte e come pensi di gestire le posizioni di 'O' e 'A' nella tua griglia?"} \\
\\ 
\\ 
\texttt{Qual è il tuo prossimo comando?} \\
            }
        }
    }
    & \\ \\

    \theutterance \stepcounter{utterance}  
    & & & \multicolumn{4}{p{0.6\linewidth}}{
        \cellcolor[rgb]{0.9,0.9,0.9}{
            \makecell[{{p{\linewidth}}}]{
                \texttt{\tiny{[GM$\langle$P2]}}
                \texttt{COMUNICA: Sono d'accordo con le tue proposte di spostare 'O' in (6,4) e 'A' in (1,5). Nella mia griglia, 'O' è ancora alla riga 2, colonna 2, e 'A' è alla riga 7, colonna 5. Potrei spostare 'O' in (6,4) e 'A' in (1,5) per allinearli con gli altri oggetti. Tuttavia, vorrei confermare che non ci siano conflitti con le tue mosse e che le posizioni finali siano coerenti con il nostro obiettivo. Prima di eseguire queste mosse, vorrei chiederti se pensi che siamo pronti a terminare il gioco dopo aver spostato 'O' e 'A'. Se si, potremmo entrambi inviare il messaggio "finito" per concludere il gioco. Quale è la tua opinione al riguardo?} \\
            }
        }
    }
     \\ \\

    \theutterance \stepcounter{utterance}  
    & & & \multicolumn{2}{p{0.3\linewidth}}{
        \cellcolor[rgb]{0.9,0.9,0.9}{
            \makecell[{{p{\linewidth}}}]{
                \texttt{\tiny{[GM$|$GM]}}
                \texttt{COMUNICA: Sono d'accordo con le tue proposte di spostare 'O' in (6,4) e 'A' in (1,5). Nella mia griglia, 'O' è ancora alla riga 2, colonna 2, e 'A' è alla riga 7, colonna 5. Potrei spostare 'O' in (6,4) e 'A' in (1,5) per allinearli con gli altri oggetti. Tuttavia, vorrei confermare che non ci siano conflitti con le tue mosse e che le posizioni finali siano coerenti con il nostro obiettivo. Prima di eseguire queste mosse, vorrei chiederti se pensi che siamo pronti a terminare il gioco dopo aver spostato 'O' e 'A'. Se si, potremmo entrambi inviare il messaggio "finito" per concludere il gioco. Quale è la tua opinione al riguardo?} \\
            }
        }
    }
    & & \\ \\

    \theutterance \stepcounter{utterance}  
    & & \multicolumn{4}{p{0.6\linewidth}}{
        \cellcolor[rgb]{0.9,0.9,0.9}{
            \makecell[{{p{\linewidth}}}]{
                \texttt{\tiny{[P1$\langle$GM]}}
                \texttt{Il tuo messaggio è stato inoltrato all'altro giocatore.} \\
\\ 
\\ 
\texttt{You are currently playing round 26 of maximum 28.} \\
\texttt{You have collectively accumulated 6 of 16 penalties.} \\
\texttt{L'altro giocatore ha inviato questo messaggio:} \\
\\ 
\\ 
\texttt{"Sono d'accordo con le tue proposte di spostare 'O' in (6,4) e 'A' in (1,5). Nella mia griglia, 'O' è ancora alla riga 2, colonna 2, e 'A' è alla riga 7, colonna 5. Potrei spostare 'O' in (6,4) e 'A' in (1,5) per allinearli con gli altri oggetti. Tuttavia, vorrei confermare che non ci siano conflitti con le tue mosse e che le posizioni finali siano coerenti con il nostro obiettivo. Prima di eseguire queste mosse, vorrei chiederti se pensi che siamo pronti a terminare il gioco dopo aver spostato 'O' e 'A'. Se si, potremmo entrambi inviare il messaggio "finito" per concludere il gioco. Quale è la tua opinione al riguardo?"} \\
\\ 
\\ 
\texttt{Qual è il tuo prossimo comando?} \\
            }
        }
    }
    & \\ \\

    \theutterance \stepcounter{utterance}  
    & \multicolumn{4}{p{0.6\linewidth}}{
        \cellcolor[rgb]{0.9,0.9,0.9}{
            \makecell[{{p{\linewidth}}}]{
                \texttt{\tiny{[P1$\rangle$GM]}}
                \texttt{COMUNICA: Sono d'accordo che le posizioni finali di 'O' in (6,4) e 'A' in (1,5) siano coerenti con il nostro obiettivo, e credo che siamo pronti a terminare il gioco dopo aver spostato 'O' e 'A', quindi sono pronto a inviare il messaggio "finito" per concludere il gioco, ma vorrei aspettare che tu abbia confermato di essere pronto a fare lo stesso, quindi ti chiedo: sei pronto a terminare il gioco? Se si, rispondi con "finito?" e poi invieremo entrambi il messaggio "finito" per concludere il gioco.} \\
            }
        }
    }
    & & \\ \\

    \theutterance \stepcounter{utterance}  
    & & & \multicolumn{2}{p{0.3\linewidth}}{
        \cellcolor[rgb]{0.9,0.9,0.9}{
            \makecell[{{p{\linewidth}}}]{
                \texttt{\tiny{[GM$|$GM]}}
                \texttt{COMUNICA: Sono d'accordo che le posizioni finali di 'O' in (6,4) e 'A' in (1,5) siano coerenti con il nostro obiettivo, e credo che siamo pronti a terminare il gioco dopo aver spostato 'O' e 'A', quindi sono pronto a inviare il messaggio "finito" per concludere il gioco, ma vorrei aspettare che tu abbia confermato di essere pronto a fare lo stesso, quindi ti chiedo: sei pronto a terminare il gioco? Se si, rispondi con "finito?" e poi invieremo entrambi il messaggio "finito" per concludere il gioco.} \\
            }
        }
    }
    & & \\ \\

    \theutterance \stepcounter{utterance}  
    & & \multicolumn{4}{p{0.6\linewidth}}{
        \cellcolor[rgb]{0.9,0.9,0.9}{
            \makecell[{{p{\linewidth}}}]{
                \texttt{\tiny{[GM$\rangle$P2]}}
                \texttt{Il tuo messaggio è stato inoltrato all'altro giocatore.} \\
\\ 
\\ 
\texttt{You are currently playing round 27 of maximum 28.} \\
\texttt{You have collectively accumulated 6 of 16 penalties.} \\
\texttt{L'altro giocatore ha inviato questo messaggio:} \\
\\ 
\\ 
\texttt{"Sono d'accordo che le posizioni finali di 'O' in (6,4) e 'A' in (1,5) siano coerenti con il nostro obiettivo, e credo che siamo pronti a terminare il gioco dopo aver spostato 'O' e 'A', quindi sono pronto a inviare il messaggio "finito" per concludere il gioco, ma vorrei aspettare che tu abbia confermato di essere pronto a fare lo stesso, quindi ti chiedo: sei pronto a terminare il gioco? Se si, rispondi con "finito?" e poi invieremo entrambi il messaggio "finito" per concludere il gioco."} \\
\\ 
\\ 
\texttt{Qual è il tuo prossimo comando?} \\
            }
        }
    }
    & \\ \\

    \theutterance \stepcounter{utterance}  
    & & & \multicolumn{4}{p{0.6\linewidth}}{
        \cellcolor[rgb]{0.9,0.9,0.9}{
            \makecell[{{p{\linewidth}}}]{
                \texttt{\tiny{[GM$\langle$P2]}}
                \texttt{COMUNICA: finito?} \\
            }
        }
    }
     \\ \\

    \theutterance \stepcounter{utterance}  
    & & & \multicolumn{2}{p{0.3\linewidth}}{
        \cellcolor[rgb]{0.9,0.9,0.9}{
            \makecell[{{p{\linewidth}}}]{
                \texttt{\tiny{[GM$|$GM]}}
                \texttt{COMUNICA: finito?} \\
            }
        }
    }
    & & \\ \\

    \theutterance \stepcounter{utterance}  
    & & \multicolumn{4}{p{0.6\linewidth}}{
        \cellcolor[rgb]{0.9,0.9,0.9}{
            \makecell[{{p{\linewidth}}}]{
                \texttt{\tiny{[P1$\langle$GM]}}
                \texttt{Il tuo messaggio è stato inoltrato all'altro giocatore.} \\
\\ 
\\ 
\texttt{You are currently playing round 27 of maximum 28.} \\
\texttt{You have collectively accumulated 6 of 16 penalties.} \\
\texttt{L'altro giocatore ha inviato questo messaggio:} \\
\\ 
\\ 
\texttt{"finito?"} \\
\\ 
\\ 
\texttt{Qual è il tuo prossimo comando?} \\
            }
        }
    }
    & \\ \\

    \theutterance \stepcounter{utterance}  
    & \multicolumn{4}{p{0.6\linewidth}}{
        \cellcolor[rgb]{0.9,0.9,0.9}{
            \makecell[{{p{\linewidth}}}]{
                \texttt{\tiny{[P1$\rangle$GM]}}
                \texttt{COMUNICA: finito!} \\
            }
        }
    }
    & & \\ \\

    \theutterance \stepcounter{utterance}  
    & & & \multicolumn{2}{p{0.3\linewidth}}{
        \cellcolor[rgb]{0.9,0.9,0.9}{
            \makecell[{{p{\linewidth}}}]{
                \texttt{\tiny{[GM$|$GM]}}
                \texttt{COMUNICA: finito!} \\
            }
        }
    }
    & & \\ \\

    \theutterance \stepcounter{utterance}  
    & & & \multicolumn{2}{p{0.3\linewidth}}{
        \cellcolor[rgb]{0.9,0.9,0.9}{
            \makecell[{{p{\linewidth}}}]{
                \texttt{\tiny{[GM$|$GM]}}
                \texttt{true} \\
            }
        }
    }
    & & \\ \\

    \theutterance \stepcounter{utterance}  
    & & & \multicolumn{2}{p{0.3\linewidth}}{
        \cellcolor[rgb]{0.9,0.9,0.9}{
            \makecell[{{p{\linewidth}}}]{
                \texttt{\tiny{[GM$|$GM]}}
                \texttt{* success: True} \\
\texttt{* lose: False} \\
\texttt{* aborted: False} \\
\texttt{{-}{-}{-}{-}{-}{-}{-}} \\
\texttt{* Shifts: 4.00} \\
\texttt{* Max Shifts: 12.00} \\
\texttt{* Min Shifts: 6.00} \\
\texttt{* End Distance Sum: 12.23} \\
\texttt{* Init Distance Sum: 30.19} \\
\texttt{* Expected Distance Sum: 29.33} \\
\texttt{* Penalties: 6.00} \\
\texttt{* Max Penalties: 16.00} \\
\texttt{* Rounds: 27.00} \\
\texttt{* Max Rounds: 28.00} \\
\texttt{* Object Count: 7.00} \\
            }
        }
    }
    & & \\ \\

    \theutterance \stepcounter{utterance}  
    & & & \multicolumn{2}{p{0.3\linewidth}}{
        \cellcolor[rgb]{0.9,0.9,0.9}{
            \makecell[{{p{\linewidth}}}]{
                \texttt{\tiny{[GM$|$GM]}}
                \texttt{Initial Grids:} \\
\texttt{```} \\
\texttt{Player 1:     Player 2:} \\
\texttt{ 1234567       1234567} \\
\texttt{╔═══╤═╤═╗     ╔═══╤═╤═╗} \\
\texttt{║◌◌◌│◌│◌║ 1   ║◌MS│◌│◌║ 1} \\
\texttt{║◌◌M│◌├─╢ 2   ║◌O◌│◌├─╢ 2} \\
\texttt{║AT◌├─┘◌║ 3   ║◌◌◌├─┘◌║ 3} \\
\texttt{║◌◌P│◌◌◌║ 4   ║◌◌◌│◌◌◌║ 4} \\
\texttt{╟─┬─┼───╢ 5   ╟─┬─┼───╢ 5} \\
\texttt{║◌│◌│S◌O║ 6   ║◌│◌│◌◌D║ 6} \\
\texttt{║◌│◌│◌D◌║ 7   ║T│◌│◌PA║ 7} \\
\texttt{╚═╧═╧═══╝     ╚═╧═╧═══╝} \\
\texttt{```} \\
\\ 
\texttt{Final Grids:} \\
\texttt{```} \\
\texttt{Player 1:     Player 2:} \\
\texttt{ 1234567       1234567} \\
\texttt{╔═══╤═╤═╗     ╔═══╤═╤═╗} \\
\texttt{║P◌◌│◌│◌║ 1   ║PM◌│◌│◌║ 1} \\
\texttt{║◌◌M│S├─╢ 2   ║◌O◌│S├─╢ 2} \\
\texttt{║A◌◌├─┘◌║ 3   ║◌◌◌├─┘◌║ 3} \\
\texttt{║◌T◌│O◌D║ 4   ║◌T◌│◌◌D║ 4} \\
\texttt{╟─┬─┼───╢ 5   ╟─┬─┼───╢ 5} \\
\texttt{║◌│◌│◌◌◌║ 6   ║◌│◌│◌◌◌║ 6} \\
\texttt{║◌│◌│◌◌◌║ 7   ║◌│◌│◌◌A║ 7} \\
\texttt{╚═╧═╧═══╝     ╚═╧═╧═══╝} \\
\texttt{```} \\
            }
        }
    }
    & & \\ \\

\end{supertabular}
}

\end{document}
